%% 
%% Copyright 2007-2024 Elsevier Ltd
%% 
%% This file is part of the 'Elsarticle Bundle'.
%% ---------------------------------------------
%% 
%% It may be distributed under the conditions of the LaTeX Project Public
%% License, either version 1.3 of this license or (at your option) any
%% later version.  The latest version of this license is in
%%    http://www.latex-project.org/lppl.txt
%% and version 1.3 or later is part of all distributions of LaTeX
%% version 1999/12/01 or later.
%% 
%% The list of all files belonging to the 'Elsarticle Bundle' is
%% given in the file `manifest.txt'.
%% 
%% Template article for Elsevier's document class `elsarticle'
%% with numbered style bibliographic references
%% SP 2008/03/01
%% $Id: elsarticle-template-num.tex 249 2024-04-06 10:51:24Z rishi $
%%
%\documentclass[preprint,12pt]{elsarticle}
\documentclass[preprint,12pt]{elsarticle}
%\documentclass[preprint,authoryear,12pt, times,3p]{elsarticle}
\biboptions{authoryear}

%% Use the option review to obtain double line spacing
% \documentclass[authoryear,preprint,review,12pt]{elsarticle}

%% Use the options 1p,twocolumn; 3p; 3p,twocolumn; 5p; or 5p,twocolumn
%% for a journal layout:
% \documentclass[final,1p,times]{elsarticle}
%% \documentclass[final,1p,times,twocolumn]{elsarticle}
% \documentclass[final,3p,times]{elsarticle}
%% \documentclass[final,3p,times,twocolumn]{elsarticle}
% \documentclass[final,5p,times]{elsarticle}
%% \documentclass[final,5p,times,twocolumn]{elsarticle}

%% For including figures, graphicx.sty has been loaded in
%% elsarticle.cls. If you prefer to use the old commands
%% please give \usepackage{epsfig}

%% The amssymb package provides various useful mathematical symbols
\usepackage{amssymb}
%% The amsmath package provides various useful equation environments.
\usepackage{amsmath}
%% The amsthm package provides extended theorem environments
%% \usepackage{amsthm}
\usepackage{algorithm}
\usepackage{algpseudocode} % from algorithmicx
\usepackage{booktabs}
\usepackage{graphicx} % para \resizebox
\usepackage{threeparttable} % para notas bajo la tabla (opcional pero recomendado)

\usepackage{url} % habilita \url{...}
% opcional (recomendado):
\usepackage[hidelinks]{hyperref}

\usepackage{comment}
\setlength{\emergencystretch}{2em}




% ---------------- Theorems ----------------
\newtheorem{definition}{Definition}
\newtheorem{assumption}{Assumption}
\newtheorem{proposition}{Proposition}

% ---------------- Macros ----------------
\newcommand{\E}{\mathbb{E}}
\newcommand{\Prob}{\mathbb{P}}
\newcommand{\CVaR}{\mathrm{CVaR}}
\newcommand{\VaR}{\mathrm{VaR}}
\newcommand{\R}{\mathbb{R}}
\newcommand{\1}{\mathbf{1}}


%% The lineno packages adds line numbers. Start line numbering with
%% \begin{linenumbers}, end it with \end{linenumbers}. Or switch it on
%% for the whole article with \linenumbers.
%% \usepackage{lineno}



\journal{Transportation Research Interdisciplinary Perspectives}

\begin{document}

\begin{frontmatter}

%% Title, authors and addresses

%% use the tnoteref command within \title for footnotes;
%% use the tnotetext command for theassociated footnote;
%% use the fnref command within \author or \affiliation for footnotes;
%% use the fntext command for theassociated footnote;
%% use the corref command within \author for corresponding author footnotes;
%% use the cortext command for theassociated footnote;
%% use the ead command for the email address,
%% and the form \ead[url] for the home page:
%% \title{Title\tnoteref{label1}}
%% \tnotetext[label1]{}
%% \author{Name\corref{cor1}\fnref{label2}}
%% \ead{email address}
%% \ead[url]{home page}
%% \fntext[label2]{}
%% \cortext[cor1]{}
%% \affiliation{organization={},
%%             addressline={},
%%             city={},
%%             postcode={},
%%             state={},
%%             country={}}
%% \fntext[label3]{}


\title{Risk-Aware Last-Mile Routing under Persistent Disruptions:
	Tail-Risk Planning and Selective Dynamic Recourse}



%CVaR-Guided Reinforcement Learning for Reliable Vehicle Routing under Realistic Travel-Time Distributions
%% use optional labels to link authors explicitly to addresses:
%% \author[label1,label2]{}
%% \affiliation[label1]{organization={},
%%             addressline={},
%%             city={},
%%             postcode={},
%%             state={},
%%             country={}}
%%
%% \affiliation[label2]{organization={},
%%             addressline={},
%%             city={},
%%             postcode={},
%%             state={},
%%             country={}}

%\author{} %% Author name

%% Author affiliation
%\affiliation{organization={},%Department and Organization
 %           addressline={}, 
 %           city={},
 %           postcode={}, 
 %           state={},
 %           country={}}


%% Author and affiliation
\author[1]{Mauricio Herrera-Mar\'in\corref{cor1}}
\ead{mherrera@udd.cl}

\cortext[cor1]{Corresponding author.}

\affiliation[1]{%
	organization={Faculty of Engineering, Universidad del Desarrollo},
	addressline={Avd. Plaza 680},
	city={Santiago},
	postcode={7550000},
	country={Chile}
}




\begin{abstract}
	
	Last-mile delivery systems are increasingly exposed to disruptions that are
	temporally persistent rather than independent across time. We study a stochastic
	capacitated vehicle routing problem with time windows (CVRPTW) in which an
	exogenous Hawkes process represents clustered disruptions, while route-dependent
	schedule propagation determines how those disruptions translate into service
	delay. The framework combines risk-aware Adaptive Large Neighbourhood Search
	(ALNS), learning-augmented operator selection, and causal Dynamic Data-Driven
	Application Systems (DDDAS) recourse. Experiments use nine Athens-derived
	operational instances under controlled stress scenarios, with common random
	numbers and matched stochastic search budgets.
	
	Increasing disruption persistence substantially raises upper-tail tardiness
	even when the continuous-time mean environmental travel-time multiplier is held
	approximately constant: for an unchanged benchmark route,
	$\CVaR_{0.95}$ tardiness increases from $60.74$ to $113.24$ between the
	low- and high-persistence regimes. The strongest algorithmic result is obtained
	from explicit tail-risk optimization. Across all nine operational days,
	CVaR-oriented ALNS improves the aligned tail-risk objective relative to
	mean-oriented ALNS by $8.61$ on average and reduces
	$\CVaR_{0.95}$ tardiness by $7.28$, with positive paired gains on all nine
	days. By contrast, feed-forward and recurrent reinforcement-learning operator
	selectors show no robust advantage over adaptive ALNS or a uniform-selection
	control under matched online budgets and zero-shot transfer.
	
	Sequential DDDAS recourse is selective rather than uniformly beneficial.
	Among 523 recourse attempts, 87 produce route modifications; all improve the
	local conditional risk objective at decision time, yet realized composite cost
	improves in only 18 of the 65 runs containing a route change. These results
	distinguish persistent environmental risk, ex-ante tail-aware planning, and
	ex-post adaptive recourse, and indicate that the most reproducible gains arise
	from explicit CVaR-based preoperational routing rather than from additional
	learning or continual rerouting.
	
\end{abstract}

















%%Graphical abstract
%\begin{graphicalabstract}
%\includegraphics{grabs}
%\end{graphicalabstract}

%%Research highlights
%\begin{highlights}
%	\item We develop a unified OR–ML framework for robust VRPTW integrating OR-Tools, ALNS, reinforcement learning and CVaR optimisation.
%	\item Realistic travel-time uncertainty is modelled using PERT and triangular distributions calibrated from multi-scenario field data.
%	\item Risk-neutral RL--ALNS improves expected performance but remains vulnerable to congestion spikes and heavy-tailed uncertainty.
%	\item The proposed RL--CVaR method reduces 95\%-CVaR of tardiness by one to two orders of magnitude with minimal increase in travel distance.
%	\item Results demonstrate that risk-sensitive operator learning yields reliable and interpretable routing policies for last-mile logistics.
%\end{highlights}

%% Keywords
%\begin{keyword}
%% keywords here, in the form: keyword \sep keyword

%% PACS codes here, in the form: \PACS code \sep code

%% MSC codes here, in the form: \MSC code \sep code
%% or \MSC[2008] code \sep code (2000 is the default)

%\end{keyword}


\begin{keyword}
	Last-mile logistics \sep
	CVRPTW \sep
	Persistent disruptions \sep
	Conditional Value-at-Risk \sep
	Hawkes processes \sep
	Adaptive large neighbourhood search \sep
	Dynamic routing
\end{keyword}

\end{frontmatter}

%% Add \usepackage{lineno} before \begin{document} and uncomment 
%% following line to enable line numbers
%% \linenumbers

%% main text
%%

\section{Introduction}

Urban last-mile logistics increasingly operate in regimes where service reliability
is dominated not by average conditions, but by rare yet persistent disruptions that
propagate through tightly coupled delivery schedules.
Congestion waves, local incidents, and the depletion of schedule slack interact over time,
producing delay cascades whose impact may far exceed that predicted by classical
memoryless uncertainty models.
Travel-time conditions in urban routing can vary systematically over time, and recent
reviews of time-dependent and stochastic dynamic routing emphasize that temporal structure,
information arrival, and real-time reoptimization materially affect route quality
\citep{Ritzinger2016, ZhangVanWoensel2023, Adamo2024, Mardesic2024}.
In such environments, routing plans that are optimal---or even robust---under static or
weakly stochastic assumptions may exhibit substantial fragility once executed
sequentially over time.

The Capacitated Vehicle Routing Problem with Time Windows (CVRPTW) is a cornerstone of
transportation science and operations research, with more than five decades of
methodological development
\citep{Solomon1987, TothVigo2002, Laporte2009}.
State-of-the-art exact, heuristic, and metaheuristic approaches---most notably Large
Neighbourhood Search (LNS) and Adaptive Large Neighbourhood Search (ALNS)---are capable of
producing high-quality solutions for large-scale deterministic instances
\citep{BraysyGendreau2005, Ropke2006ALNS, Christiaens2020SurveyVRPTW}.
Modern industrial solvers such as Google OR-Tools further illustrate the practical
scalability of deterministic vehicle-routing optimization under nominal travel-time
assumptions \citep{ortools2023}.
In their standard deterministic use, however, such formulations do not explicitly
represent temporal persistence in the disruption process or the resulting exposure of a
route to clustered periods of degraded travel conditions.

A substantial literature has demonstrated that ignoring travel-time uncertainty can lead
to considerable degradation of operational performance.
Stochastic, robust, and distributionally robust routing formulations address this problem
through scenario-based, probabilistic, or ambiguity-set representations of uncertainty
\citep{Ritzinger2016, ZhangEtAl2021DRODeadlines, Mardesic2024}.
Risk-sensitive criteria provide an additional mechanism for controlling service
unreliability.
In particular, Conditional Value-at-Risk (CVaR) offers a coherent way to target the upper
tail of tardiness rather than optimizing expected performance alone
\citep{RockafellarUryasev2000}.
Yet temporal dependence and persistence remain less developed than marginal travel-time
uncertainty itself: scenario-based formulations can represent heterogeneous and even
correlated outcomes without necessarily describing how disruption episodes cluster,
decay, and become observable during sequential route execution.

In parallel, learning-augmented optimization has gained prominence in vehicle routing.
Reinforcement learning (RL) has been applied to end-to-end route construction
\citep{Nazari2018RL, Kool2019Attention}, while recent surveys document a broader shift toward
learning-augmented routing and hybrid combinations of learned policies with classical
optimization or improvement heuristics
\citep{Bogyrbayeva2024MLVRP, ZhouEtAl2025LearningRouting, WangEtAl2025RLALNS}.
These hybrid approaches are attractive because learning can guide expensive search
decisions while combinatorial operators preserve the structural constraints of the
routing problem.
However, an important question remains insufficiently resolved: whether learned operator
selection provides an operational advantage once search effort, stochastic evaluation
budget, offline training cost, and generalization to unseen instances are controlled.
This question becomes particularly relevant when the optimization criterion is not mean
performance but tail risk, and when the state available to a routing policy evolves with
the observed disruption history.

From a systems perspective, two distinct temporal mechanisms must be separated.
First, disruption conditions may themselves exhibit \emph{exogenous persistence}: recent
events contain information about the likelihood of further events even though routing
decisions do not alter the disruption-generating process.
Second, delays propagate \emph{through the delivery schedule}: a slow arc changes later
departure times, consumes time-window slack, and can transform subsequent travel-time
realizations into tardiness.
The second mechanism can therefore amplify the operational consequences of the first
without generating feedback from the route to the disruption process.
This distinction is important because a route may be exposed to the same external
disruption environment as another route while experiencing different consequences simply
because its arc departures occur at different times.

We represent the first mechanism using an exogenous self-exciting Hawkes process
\citep{hawkes1971, bacry2015hawkes}.
Recent disruption events generate a decaying excitation state, which is mapped to an
\emph{Operational Memory Index} (OMI) and then to a multiplicative travel-time effect.
The Hawkes process is common to the operating environment and is not affected by realized
routing delays.
For the exponential kernel used here, the excitation admits a finite-dimensional Markov
representation; accordingly, the system is history-dependent relative to a conventional
routing state that omits this excitation variable, but it becomes Markovian once the
environmental memory state is included.
Route dependence then arises because different routing decisions induce different
departure times and hence different exposure to the same evolving OMI trajectory.

Against this background, the present study asks three related but distinct questions.
First, how does disruption persistence affect mean and tail routing risk when differences
in average environmental impact are controlled?
Second, does explicitly optimizing tail risk through CVaR provide a reproducible advantage
over mean-risk optimization under a common stochastic evaluation protocol?
Third, once a risk-aware preoperational route is available, how much additional value is
provided by causal online recourse based only on information observed up to each decision
epoch?
A complementary learning experiment asks whether feed-forward or recurrent RL-based
operator selection improves upon the underlying adaptive search once computational
budgets and generalization are controlled.

To address these questions, we formulate a stochastic CVRPTW with two-dimensional vehicle
capacity and a common risk-sensitive objective,
\[
J_{\rho}(x)
=
\lambda_D D(x)
+
\lambda_T \rho(Z_x)
+
\lambda_W O_W(x)
+
\lambda_V O_V(x),
\]
where $D(x)$ denotes route distance, $Z_x$ is cumulative tardiness under stochastic
execution, $O_W(x)$ and $O_V(x)$ denote weight and volume capacity violations, and
$\rho$ is either the mean or $\operatorname{CVaR}_{\alpha}$.
The same objective structure is used across the principal search and learning
comparisons.
We implement a genuine adaptive ALNS and compare mean-oriented and CVaR-oriented search
under common random numbers and matched candidate--scenario evaluation budgets.
Feed-forward and gated recurrent unit (GRU) policies are evaluated as learned
destroy--repair operator selectors rather than as replacements for the combinatorial
solver.

The framework is then extended to sequential operation through a Dynamic Data-Driven
Application Systems (DDDAS) recourse layer.
At each eligible decision epoch, only the disruption events and travel-time information
observed before departure are available.
When the OMI enters a prospectively specified high-memory episode, the system may
reconsider the residual route of the affected vehicle.
Future disruption events are generated conditionally from the observed Hawkes state,
already traversed arcs retain their observed travel times, and unobserved travel times
remain stochastic.
The resulting recourse problem is therefore causal and local to the residual vehicle
suffix.
Online search is explicitly bounded, so a disruption trigger should be interpreted as a
request to \emph{reconsider} the incumbent route rather than as an instruction to reroute.

The empirical analysis uses multi-day last-mile delivery instances derived from real
operations in Athens, Greece \citep{Vrani2025DataInBrief}.
The operational data provide the customer, demand, spatial, and travel-time structure of
the routing instances, while moderately tightened time windows and the Hawkes disruption
process define controlled stress scenarios used to expose reliability differences among
methods.
The disruption process should therefore be interpreted as an experimental persistence
model rather than as a Hawkes process calibrated from observed Athens incident arrivals.
Nine operational days are used for the principal cross-instance evaluation, and three
prospectively selected representative days are examined more deeply for sequential
DDDAS behavior.

The experiments reveal a differentiated picture.
Increasing disruption persistence substantially increases tail tardiness risk even under
an auxiliary control that matches the continuous-time mean environmental travel-time
multiplier across persistence regimes.
At the planning stage, CVaR-oriented ALNS consistently improves the tail-oriented
objective relative to mean-oriented ALNS across all nine operational days, indicating
that explicit tail-risk optimization provides a reproducible benefit beyond optimizing
expected tardiness alone.
By contrast, neither feed-forward nor recurrent learned operator selection exhibits a
robust advantage over the adaptive search or a uniform operator-selection control once
search budgets, validation, and transfer to other days are considered.
Thus, recurrent state augmentation is not assumed to be beneficial merely because the
environment contains temporal persistence.

The sequential DDDAS experiments provide a further distinction between
\emph{decision-time conditional risk} and \emph{realized ex-post performance}.
Across the confirmatory sequential experiment, recourse was triggered frequently but
actual route modification was selective.
Every executed route modification improved the local conditional search objective relative
to its incumbent residual route, yet these ex-ante conditional improvements did not
translate systematically into lower realized composite cost on the subsequently observed
trajectory.
The online mechanism is therefore best interpreted as bounded, state-dependent recourse
under uncertainty rather than as a guaranteed realized-cost improvement procedure.
This result also illustrates why evaluating an online risk-sensitive controller solely by
the realized outcome of a single stochastic trajectory can obscure the distinction between
the objective optimized when the decision is made and the realization that subsequently
occurs.

The contributions of this work are fourfold.
\begin{enumerate}
	\item We develop a persistent-disruption formulation of the stochastic CVRPTW that
	explicitly separates an exogenous Hawkes disruption process from route-dependent
	schedule propagation and exposure.
	
	\item We provide a controlled assessment of temporal persistence and tail risk,
	including a mean-impact-matched experiment that isolates changes in disruption
	clustering from changes in average environmental travel-time impact.
	
	\item We establish a common-budget comparison of mean-oriented ALNS, CVaR-oriented
	ALNS, and learning-augmented operator-selection policies, including feed-forward and
	recurrent alternatives and zero-shot transfer across operational days.
	
	\item We implement and evaluate a causal, computationally bounded DDDAS recourse
	mechanism in which decisions depend only on information available at the decision
	epoch, thereby quantifying both the potential and the limitations of online
	risk-sensitive adaptation.
\end{enumerate}

Taken together, the results suggest that resilience under persistent disruption should not
be identified with a single algorithmic mechanism.
Explicit tail-risk optimization provides the clearest and most reproducible improvement
at the preoperational planning stage, whereas learning-based search control and online
recourse provide additional adaptive capabilities whose value depends on the information,
computational budget, and realized trajectory.
This separation between environmental persistence, risk-aware planning, and sequential
recourse provides a more transparent basis for designing data-driven decision support for
disruption-prone last-mile logistics.




% ============================================================
\section{Related Work}
\label{sec:related_work}

\subsection{Vehicle routing with time windows under uncertainty}

The vehicle routing problem with time windows (VRPTW) has developed from the
classical deterministic formulation of \citet{Solomon1987} into a broad family
of stochastic, robust, time-dependent, and dynamic routing models
\citep{TothVigo2002, Laporte2009, BraysyGendreau2005}.
Recent reviews stress that uncertainty is not characterized only by its marginal
magnitude: the timing of information, temporal variation of travel conditions,
and the possibility of reoptimization are central modelling choices
\citep{Ritzinger2016, ZhangVanWoensel2023, Mardesic2024, Adamo2024}.
Our setting keeps the daily customer set fixed before dispatch, while a common
exogenous disruption environment evolves during execution. Dynamic adaptation
therefore enters through causal residual-route reconsideration, not through
new customer arrivals.

\subsection{Risk-averse routing and tail performance metrics}

Expectation-based objectives can understate operational exposure when a small
number of severe delays dominate service penalties. Risk-averse routing therefore
uses chance constraints, robust or distributionally robust formulations, quantiles,
and coherent tail-risk measures such as Conditional Value-at-Risk (CVaR)
\citep{RockafellarUryasev2000}. Distributionally robust vehicle-routing models have,
for example, used CVaR approximations to control deadline violations when the
underlying distribution is uncertain \citep{ZhangEtAl2021DRODeadlines}.
The contribution here is not CVaR itself, but a controlled test of whether explicit
CVaR optimization remains beneficial when uncertainty is modulated by a persistent
latent disruption state and when average environmental impact is held fixed.

\subsection{Metaheuristics and adaptive large neighbourhood search}

Adaptive Large Neighbourhood Search (ALNS) is a well-established framework for
rich routing problems \citep{Ropke2006ALNS}. Its defining feature is adaptive
operator selection based on observed search performance rather than a fixed
neighbourhood schedule. A recent systematic review of 211 VRP studies provides a
unified taxonomy of removal and insertion operators and reports strong performance
for foresight-based insertions such as regret insertion \citep{Voigt2025ALNS}.
ALNS also remains relevant in explicitly dynamic VRPTW variants; for example,
\citet{WangSunHuang2024} develop an ALNS reoptimization framework for a multi-depot
dynamic VRPTW. These results motivate using a genuinely adaptive ALNS as the main
combinatorial benchmark here. Mean- and CVaR-oriented variants share the same
adaptive machinery, neighbourhoods, and stochastic budget, isolating the optimized
risk functional.

\subsection{Learning-augmented routing and reinforcement learning}

Learned methods now span end-to-end route construction and augmentation of
classical search. Early policy-gradient and attention models demonstrated that
neural policies can construct competitive routing solutions
\citep{Nazari2018RL, Kool2019Attention}. A 2024 survey synthesizes machine-learning
approaches to VRPs \citep{Bogyrbayeva2024MLVRP}, and a 2025 review organizes recent
routing work into construction- and improvement-based learning paradigms
\citep{ZhouEtAl2025LearningRouting}. Hybrid RL--ALNS methods continue to appear,
including PPO-guided operator control for VRPTW \citep{WangEtAl2025RLALNS}.
The present learning component follows this hybrid paradigm: feed-forward and
recurrent policies select destroy operators, while feasibility, repair, stochastic
scoring, and acceptance remain in the common ALNS engine. This permits an
incremental test against both adaptive ALNS and uniform operator selection, with
offline training cost reported separately from the matched online search budget.

\subsection{Temporal dependence, self-excitation, and dynamic control}

Time-dependent and stochastic dynamic routing already establish that route quality
can depend strongly on when travel conditions change and when new information is
available \citep{ZhangVanWoensel2023, Adamo2024}. Self-exciting point processes
provide a complementary mechanism for representing event clustering. Hawkes'
original construction \citep{hawkes1971} has since been used broadly for clustered
event processes, and a recent transportation example models road accidents in Rome
with a spatio-temporal Hawkes process \citep{AlaimoDiLoroEtAl2025}.
Here the Hawkes process is deliberately an \emph{experimental exogenous disruption
model}: it is not calibrated from Athens incident arrivals and routing decisions do
not feed back into its intensity. Its excitation state determines a common
Operational Memory Index (OMI), while different routes experience different exposure
because they depart onto arcs at different times.

Dynamic Data-Driven Application Systems (DDDAS) provide a natural systems concept
for the execution stage because observations acquired during operation can update
the state used for subsequent decisions \citep{Darema2004DDDAS}. In our
implementation, this principle is causal and local: only events and completed travel
times observed before a decision epoch condition the residual problem, and only the
unexecuted suffix of the affected vehicle can be revised.

\subsection{Positioning of this work}

The gap addressed here lies at the intersection of these mature streams. We study a
CVRPTW in which (i) disruption events exhibit explicit temporal persistence through
an exogenous excitation state, (ii) persistence is separated from average
environmental impact through matched controls, (iii) mean- and CVaR-oriented routing
are compared using the same adaptive ALNS and common stochastic budgets, (iv)
learned operator selection is evaluated as an incremental augmentation rather than
assumed to dominate adaptive search, and (v) online adaptation is implemented as
causal, bounded residual-route recourse. Accordingly, we do not claim that Hawkes
processes uniquely represent urban disruption or that RL is intrinsically superior;
the experiments identify which additional mechanisms are supported under the common
confirmatory protocol.

\begin{figure*}[t]
\centering
\includegraphics[width=\textwidth]{figures/main/pdf/Fig1_architecture_v3.pdf}
\caption{\textbf{Conceptual architecture of persistent-risk routing.} The disruption process is exogenous to routing. Baseline OD uncertainty and the Hawkes-derived OMI jointly determine effective arc travel times, while route choices determine the departure times at which the common environmental state is sampled. Preoperational risk-aware planning is followed by sequential execution. The DDDAS layer uses only the observed disruption prefix and completed travel times to reconsider the unexecuted vehicle suffix; the revised suffix then returns to execution. No route-to-Hawkes feedback is assumed (crossed dashed arrow).}
\label{fig:architecture}
\end{figure*}

\section{Problem Formulation}
\label{sec:formulation}
% ============================================================

We consider a stochastic capacitated vehicle routing problem with soft time
windows in which vehicles operate under an exogenous disruption environment with
temporal persistence. The formulation separates three components that play
different roles during execution: (i) baseline origin--destination travel-time
uncertainty, (ii) an exogenous self-exciting disruption process, and
(iii) route-dependent schedule propagation. Routing decisions affect when a
vehicle is exposed to the evolving disruption state, but they do not affect the
disruption-generating process itself.

% ------------------------------------------------------------
\subsection{Network, routes, demands, and capacity}
% ------------------------------------------------------------

Let $G=(V,A)$ be a directed graph with depot $0$ and customer set
$\mathcal{C}=\{1,\ldots,n\}$, so that $V=\{0\}\cup\mathcal{C}$.
The fleet is denoted by $\mathcal{K}=\{1,\ldots,K\}$.
Each customer $i\in\mathcal{C}$ is characterized by a weight demand $w_i$,
a volume demand $v_i$, a service duration $s_i$, and a service window
$[a_i,b_i]$.

The upper endpoint $b_i$ is treated as a service target rather than as a hard
feasibility cutoff: vehicles may serve a customer after $b_i$, but the resulting
tardiness is penalized. Early arrival is permitted and induces waiting until
$a_i$.

A routing solution can be represented either by binary arc variables
$x_{ij}^k\in\{0,1\}$ or equivalently by a collection of vehicle routes
\[
x=\{r_1,\ldots,r_K\},
\qquad
r_k=(0,i_{k,1},\ldots,i_{k,n_k},0).
\]
Each customer is assigned to exactly one route.

Let $Q^W_k$ and $Q^V_k$ denote the weight and volume capacities of vehicle $k$.
The corresponding route loads are
\[
L^W_k(x)=\sum_{i\in r_k\cap\mathcal{C}} w_i,
\qquad
L^V_k(x)=\sum_{i\in r_k\cap\mathcal{C}} v_i.
\]
Capacity excesses are
\begin{align}
	O_W(x)
	&=
	\sum_{k\in\mathcal{K}}
	\left[L^W_k(x)-Q^W_k\right]_+,
	\\
	O_V(x)
	&=
	\sum_{k\in\mathcal{K}}
	\left[L^V_k(x)-Q^V_k\right]_+,
\end{align}
where $[z]_+=\max\{z,0\}$.
The implemented destroy--repair search enforces both capacity dimensions when
constructing admissible routes; consequently, these violation terms are zero for
all reported feasible solutions and serve as safeguards in the canonical scoring
function.

The deterministic routing distance is
\begin{equation}
	D(x)
	=
	\sum_{k\in\mathcal{K}}
	\sum_{\ell=0}^{n_k}
	d_{i_{k,\ell},i_{k,\ell+1}},
	\label{eq:distance}
\end{equation}
where $i_{k,0}=i_{k,n_k+1}=0$ and $d_{ij}$ is the observed OD distance.

% ------------------------------------------------------------
\subsection{Baseline travel-time uncertainty}
\label{subsec:baseline_tt}
% ------------------------------------------------------------

For each OD pair $(i,j)$, the operational data provide optimistic,
most-likely, and pessimistic travel-time estimates,
\[
t^{\mathrm{opt}}_{ij},
\qquad
t^{\mathrm{ml}}_{ij},
\qquad
t^{\mathrm{pes}}_{ij}.
\]
These triplets are retained unchanged in the source data. For stochastic
sampling only, occasional ordering inconsistencies are regularized as
\begin{equation}
	L_{ij}
	=
	\min\left\{
	t^{\mathrm{opt}}_{ij},
	t^{\mathrm{ml}}_{ij},
	t^{\mathrm{pes}}_{ij}
	\right\},
	\qquad
	M_{ij}=t^{\mathrm{ml}}_{ij},
	\qquad
	U_{ij}
	=
	\max\left\{
	t^{\mathrm{opt}}_{ij},
	t^{\mathrm{ml}}_{ij},
	t^{\mathrm{pes}}_{ij}
	\right\}.
	\label{eq:pert_regularization}
\end{equation}

The baseline stochastic travel time
$B_{ij}^{(\omega)}$ is sampled from a modified PERT distribution on
$[L_{ij},U_{ij}]$ with mode $M_{ij}$ and shape parameter $\lambda_{\mathrm{PERT}}=4$.
For $U_{ij}>L_{ij}$, the corresponding beta-shape parameters are
\begin{align}
	a_{ij}
	&=
	1+
	\lambda_{\mathrm{PERT}}
	\frac{M_{ij}-L_{ij}}{U_{ij}-L_{ij}},
	\\
	b_{ij}
	&=
	1+
	\lambda_{\mathrm{PERT}}
	\frac{U_{ij}-M_{ij}}{U_{ij}-L_{ij}}.
\end{align}
Degenerate triplets with $L_{ij}=U_{ij}$ yield a deterministic baseline time.

The regularization in Eq.~\eqref{eq:pert_regularization} is used only to define
a valid stochastic sampling distribution; it does not modify the raw OD data.
The nominal deterministic routing benchmark uses the corresponding nominal
most-likely travel-time information rather than stochastic PERT draws.

% ------------------------------------------------------------
\subsection{Exogenous persistent disruption process}
\label{subsec:hawkes}
% ------------------------------------------------------------

Temporal clustering of disruptions is represented by a univariate Hawkes
process~\citep{hawkes1971,bacry2015hawkes}. Let
$\{\tau_q\}_{q\ge1}$ denote disruption-event times. With an exponential kernel,
the conditional event intensity is
\begin{equation}
	\lambda_H(t)
	=
	\mu
	+
	\xi
	\sum_{\tau_q<t}
	\exp\!\left[-\beta(t-\tau_q)\right],
	\label{eq:hawkes_intensity}
\end{equation}
where $\mu>0$ is the baseline event rate, $\beta>0$ is the decay rate, and
$\xi>0$ is the excitation amplitude.

We parameterize the excitation amplitude as
\begin{equation}
	\xi=\eta\beta,
	\qquad 0\leq\eta<1,
	\label{eq:branching}
\end{equation}
so that $\eta=\xi/\beta$ is the Hawkes branching ratio.
For constant $\mu$, the stationary mean event intensity is
\begin{equation}
	\bar{\lambda}
	=
	\frac{\mu}{1-\eta}.
	\label{eq:hawkes_stationary_rate}
\end{equation}

Importantly, the Hawkes process is \emph{exogenous to routing}.
The event process is generated independently of route choices, realized
travel times, and realized tardiness. A delayed vehicle therefore does not
increase future Hawkes intensity. The interaction between the route and the
disruption process occurs only through the times at which the route samples
the external disruption state.

For the exponential kernel, define the excitation state
\begin{equation}
	m(t)
	=
	\xi
	\sum_{\tau_q<t}
	\exp\!\left[-\beta(t-\tau_q)\right].
	\label{eq:memory_state}
\end{equation}
Between events, $m(t)$ decays exponentially, and at a disruption event it jumps
by $\xi$. Hence the exponential Hawkes model admits a finite-dimensional Markov
embedding through $m(t)$. The process is therefore history-dependent relative
to a conventional routing state that omits this variable, but it is not
intrinsically non-Markovian once the excitation state is included.

% ------------------------------------------------------------
\subsection{Operational Memory Index and route exposure}
\label{subsec:omi}
% ------------------------------------------------------------

We summarize the instantaneous contribution of past events to the current
disruption intensity through the Operational Memory Index (OMI),
\begin{equation}
	\mathrm{OMI}(t)
	=
	\frac{m(t)}
	{\mu+m(t)}
	\in[0,1].
	\label{eq:omi}
\end{equation}
An OMI close to zero indicates that the current intensity is dominated by the
baseline component $\mu$, whereas values close to one indicate that recent
events dominate the instantaneous intensity.

OMI is an \emph{environment-level} quantity: for a fixed Hawkes realization,
the same function $\mathrm{OMI}(t)$ is available to every routing method.
Methods do not suppress or modify OMI. They can, however, experience different
\emph{route-dependent OMI exposure} because different routes depart on different
arcs at different times.

% ------------------------------------------------------------
\subsection{Effective travel times and sequential execution}
\label{subsec:execution}
% ------------------------------------------------------------

For scenario $\omega$, the effective travel time on arc $(i,j)$ when that arc
is started at time $t$ is
\begin{equation}
	T_{ij}^{(\omega)}(t)
	=
	B_{ij}^{(\omega)}
	\left[
	1+\kappa\,\mathrm{OMI}^{(\omega)}(t)
	\right],
	\label{eq:effective_tt}
\end{equation}
where $\kappa\ge0$ controls the operational impact of the disruption state.

Equation~\eqref{eq:effective_tt} defines exposure at \emph{arc departure}.
Once an arc has been started, its travel-time multiplier is fixed; events that
occur during traversal do not retrospectively modify that arc. They become
observable only at a subsequent decision epoch.

Consider route
\[
r_k=(0,i_1,\ldots,i_{n_k},0)
\]
under scenario $\omega$.
Let $A_\ell^{(\omega)}$ denote arrival time at customer $i_\ell$,
$S_\ell^{(\omega)}$ its service-start time, and
$D_\ell^{(\omega)}$ the departure time after service.
The execution recursion is
\begin{align}
	D_0^{(\omega)} &= 0,
	\\
	A_\ell^{(\omega)}
	&=
	D_{\ell-1}^{(\omega)}
	+
	T_{i_{\ell-1},i_\ell}^{(\omega)}
	\!\left(D_{\ell-1}^{(\omega)}\right),
	\\
	S_\ell^{(\omega)}
	&=
	\max\left\{
	A_\ell^{(\omega)},a_{i_\ell}
	\right\},
	\\
	D_\ell^{(\omega)}
	&=
	S_\ell^{(\omega)}
	+
	s_{i_\ell},
	\qquad
	\ell=1,\ldots,n_k.
	\label{eq:schedule_recursion}
\end{align}

Customer tardiness is measured at service start,
\begin{equation}
	Z_{i_\ell}^{(\omega)}(x)
	=
	\left[
	S_\ell^{(\omega)}-b_{i_\ell}
	\right]_+,
	\label{eq:customer_tardiness}
\end{equation}
and total scenario tardiness is
\begin{equation}
	Z_x^{(\omega)}
	=
	\sum_{k\in\mathcal{K}}
	\sum_{\ell=1}^{n_k}
	Z_{i_\ell}^{(\omega)}(x).
	\label{eq:total_tardiness}
\end{equation}

These recursions create route-level path dependence even though the disruption
process itself is exogenous. An early travel-time realization can alter later
departure times, deplete time-window slack, and thereby change the OMI values
sampled by subsequent arcs. Two methods exposed to the same Hawkes trajectory
can therefore experience different tardiness because their routes induce
different temporal exposure profiles.

% ------------------------------------------------------------
\subsection{Mean- and tail-risk objectives}
\label{subsec:risk_objective}
% ------------------------------------------------------------

For a routing solution $x$, we use the canonical risk-sensitive objective
\begin{equation}
	J_{\rho}(x)
	=
	\lambda_D D(x)
	+
	\lambda_T \rho(Z_x)
	+
	\lambda_W O_W(x)
	+
	\lambda_V O_V(x),
	\label{eq:canonical_objective}
\end{equation}
where $\rho$ is either expected total tardiness,
\[
\rho(Z_x)=\mathbb{E}[Z_x],
\]
or Conditional Value-at-Risk,
\[
\rho(Z_x)=\CVaR_{\alpha}(Z_x).
\]

All principal experiments use $\alpha=0.95$.
The same objective structure and penalty convention are used for ALNS and the
learning-augmented variants; only the risk functional $\rho$ changes between
mean-oriented and CVaR-oriented experiments.

Given a frozen scenario bank
$\Omega=\{\omega_1,\ldots,\omega_M\}$,
the sample mean is
\begin{equation}
	\widehat{\mathbb{E}}_{\Omega}[Z_x]
	=
	\frac{1}{M}
	\sum_{m=1}^{M}Z_x^{(\omega_m)}.
\end{equation}
Empirical CVaR is evaluated as the expected shortfall of the empirical
distribution. Equivalently,
\begin{equation}
	\widehat{\CVaR}_{\alpha,\Omega}(Z_x)
	=
	\min_{\tau\in\mathbb{R}}
	\left\{
	\tau+
	\frac{1}{(1-\alpha)M}
	\sum_{m=1}^{M}
	\left[
	Z_x^{(\omega_m)}-\tau
	\right]_+
	\right\}.
	\label{eq:empirical_cvar}
\end{equation}
This definition correctly handles finite samples and ties at the empirical
quantile.

Candidate and incumbent routes are compared on the \emph{same} frozen scenario
bank whenever they compete within a search step. This common-random-number
construction reduces Monte Carlo comparison noise and prevents differences
between algorithms from being driven by different stochastic realizations.

% ------------------------------------------------------------
\subsection{Optimization questions}
% ------------------------------------------------------------

The formulation leads to two distinct decision problems.

The first is a \emph{preoperational planning problem}: before execution begins,
find a route set that performs well under the full stochastic disruption
environment. We compare mean-oriented and CVaR-oriented versions of
Eq.~\eqref{eq:canonical_objective}.

The second is an \emph{online residual-routing problem}: once execution is under
way, determine whether the unexecuted suffix of an affected vehicle should be
modified using only information observed up to the current decision epoch.
This second problem is conditional, causal, and computationally bounded.
It is developed in Section~\ref{sec:methods}.


% ============================================================
\section{Methodology}
\label{sec:methods}
% ============================================================

The computational framework separates preoperational route optimization from
online recourse. The same stochastic execution model is used throughout, but
the information available to the optimizer differs between the two stages.
Preoperational methods evaluate complete routes against frozen scenario banks.
The DDDAS layer instead constructs a conditional scenario bank from the
information available at the current execution time and optimizes only the
remaining suffix of the triggered vehicle.

% ------------------------------------------------------------
\subsection{Scenario banks and common random numbers}
\label{subsec:scenario_banks}
% ------------------------------------------------------------

A scenario bank contains two stochastic components: a realization of the
exogenous Hawkes event process and a matrix of baseline OD travel-time draws.
For a fixed route and scenario, sequential execution follows
Eqs.~\eqref{eq:effective_tt}--\eqref{eq:total_tardiness}.

Separate frozen scenario banks are used for training, validation, search, and
final testing. Within a comparison, all competing methods are evaluated on the
same final test scenarios. During local search, incumbent and candidate routes
are also evaluated against the same search bank. This common-random-number
design is central to the fairness protocol because route differences are
evaluated under identical realizations of the external disruption environment.

% ------------------------------------------------------------
\subsection{Adaptive Large Neighbourhood Search}
\label{subsec:alns}
% ------------------------------------------------------------

Adaptive Large Neighbourhood Search (ALNS) provides the combinatorial search
backbone~\citep{Ropke2006ALNS}. Starting from a feasible route set $x_k$, each
iteration applies a destroy operator followed by a repair procedure to construct
a candidate $\widetilde{x}_k$.

The destroy set contains random, worst-contribution, and Shaw-related removal
operators. At each non-null proposal, the number of removed customers is
generated from a fraction uniformly sampled between $0.10$ and $0.15$ of the
current customer set, with at least two removals. Candidate reinsertion uses a
randomized regret-2 procedure with feasible-position screening under both
weight and volume capacities. High-regret choices are rank-biased rather than
selected deterministically, allowing controlled diversification while preserving
feasibility.

A candidate is scored through the stochastic objective
\[
\widehat{J}_{\rho}(x;\Omega)
\]
on the current frozen scenario bank $\Omega$. Search uses simulated-annealing
acceptance: an improving candidate is accepted, while a non-improving candidate
may be accepted according to the current temperature. The initial temperature
is $5$ and the cooling factor is $0.995$ for each evaluated non-null candidate.

The baseline ALNS is genuinely adaptive. Operator weights are initialized
equally and updated every 25 operator decisions using
\begin{equation}
	w_a^{\mathrm{new}}
	=
	0.8\,w_a^{\mathrm{old}}
	+
	0.2
	\max\left\{
	\frac{R_a}{N_a},
	10^{-6}
	\right\},
	\label{eq:alns_weight_update}
\end{equation}
where $R_a$ is accumulated operator reward and $N_a$ is its number of uses in
the adaptation window. Rewards distinguish three outcomes: creation of a global
best solution, improvement of the current solution, and accepted
non-improving movement. The corresponding rewards are $8$, $4$, and $1$,
respectively.

This adaptive baseline is used for both the mean-oriented and CVaR-oriented
search. Consequently, the comparison
\[
\text{ALNS-Mean}
\quad\text{versus}\quad
\text{ALNS-CVaR}
\]
changes the risk functional being optimized rather than replacing a
non-adaptive heuristic with an adaptive one.

% ------------------------------------------------------------
\subsection{Learning-augmented operator selection}
\label{subsec:rl}
% ------------------------------------------------------------

We additionally investigate whether learned search control can improve upon the
adaptive ALNS baseline. The learning component does not construct routes
directly. Instead, it selects the destroy neighbourhood used to propose the
next ALNS move, while the same feasibility checks, repair mechanism, stochastic
objective, and acceptance rule remain in force.

At search step $k$, the policy observes the nine-dimensional state
\begin{equation}
	s_k=
	\left(
	\bar D_k,\,
	\bar T^{\mathrm{nom}}_k,\,
	u^W_k,\,
	u^V_k,\,
	\bar Z_k,\,
	\sigma_{Z,k},\,
	p^{\mathrm{OT}}_k,\,
	\bar E^{\mathrm{OMI}}_k,\,
	E^{\mathrm{OMI,max}}_k
	\right),
	\label{eq:rl_state}
\end{equation}
where the components are, respectively, normalized route distance, normalized
nominal tardiness, maximum weight utilization, maximum volume utilization,
normalized Monte Carlo mean tardiness, normalized Monte Carlo tardiness
standard deviation, mean on-time service fraction, mean route OMI exposure,
and maximum route OMI exposure.

The last two quantities are route-exposure statistics and should not be
interpreted as changes to the environmental OMI itself. They summarize the
OMI values encountered at the route's actual arc-departure times.

Two policy classes are considered. The feed-forward policy conditions directly
on $s_k$. The recurrent policy uses a gated recurrent unit (GRU),
\begin{equation}
	h_k
	=
	\mathrm{GRU}_{\theta}(s_k,h_{k-1}),
	\qquad
	a_k
	\sim
	\pi_{\theta}(\cdot\mid h_k).
	\label{eq:gru_policy}
\end{equation}
Only the newly observed state $s_k$ is processed at step $k$ together with the
previous hidden state $h_{k-1}$; the full preceding sequence is not processed
again.

Learning is aligned with exactly the same stochastic objective used by the
search. If $x_k$ is the incumbent route set and $\widetilde{x}_k$ the proposed
candidate, both evaluated on scenario bank $\Omega_b$, the immediate learning
signal is
\begin{equation}
	r_k
	=
	\widehat{J}_{\rho}(x_k;\Omega_b)
	-
	\widehat{J}_{\rho}(\widetilde{x}_k;\Omega_b).
	\label{eq:rl_reward}
\end{equation}
Thus a positive reward corresponds to a candidate with a lower objective on the
same stochastic realizations. Simulated-annealing acceptance is handled
separately: it determines whether the candidate becomes the next search state,
but it does not redefine the reward of a rejected proposal.

Feed-forward and recurrent policies are evaluated against both the adaptive ALNS
baseline and a uniform operator-selection control. Offline policy-training cost
is accounted for separately from the online search budget, while final
method-to-method search comparisons use matched candidate--scenario evaluation
budgets. This design allows the learning component to be assessed as an
empirical augmentation rather than assuming that recurrence must be beneficial
because the environment contains temporal persistence.

% ------------------------------------------------------------
\subsection{Causal DDDAS recourse}
\label{subsec:dddas}
% ------------------------------------------------------------

The Dynamic Data-Driven Application Systems (DDDAS) component
\citep{Darema2004DDDAS} operates during route execution.
Unlike the preoperational optimizers, it has access only to information that
would have been observed by the decision time.

A decision epoch occurs immediately before a vehicle starts a new arc, after
any waiting and service at its current node have been completed. At departure
time $t_0$, the observable disruption history is
\[
\mathcal{H}_{t_0}
=
\{\tau_q:\tau_q<t_0\};
\]
events at or after $t_0$ are not available to the controller.
Similarly, baseline travel times are regarded as observed only for arcs whose
traversal has already been completed.

For the exponential Hawkes model, the observed prefix determines the current
excitation state
\begin{equation}
	s_H(t_0)
	=
	\sum_{\tau_q<t_0}
	\exp[-\beta(t_0-\tau_q)],
	\qquad
	m(t_0)=\xi s_H(t_0).
	\label{eq:conditional_hawkes_state}
\end{equation}
Future Hawkes events can therefore be generated conditionally from the pair
$(t_0,s_H(t_0))$ without access to the unobserved future realization.

Conditional scenario generation at $t_0$ follows three rules:
\begin{enumerate}
	\item the observed Hawkes prefix is held fixed and future events are sampled
	conditionally from the state in
	Eq.~\eqref{eq:conditional_hawkes_state};
	
	\item baseline travel times for already completed OD traversals retain their
	observed realized values;
	
	\item baseline travel times for unobserved OD pairs are newly sampled from
	the PERT model in Section~\ref{subsec:baseline_tt}.
\end{enumerate}
No latent future events and no unobserved entries of the originally realized OD
matrix are supplied to the replanner.

% ------------------------------------------------------------
\subsection{Memory-triggered episodes}
\label{subsec:trigger}
% ------------------------------------------------------------

Online reconsideration is governed by a prospectively specified OMI trigger.
The controller is initially armed. The first eligible departure satisfying
\begin{equation}
	\mathrm{OMI}(t_0)\ge0.55
	\label{eq:omi_trigger}
\end{equation}
opens a high-memory episode.
The controller is rearmed only after a subsequent observed departure has
$\mathrm{OMI}(t)<0.55$; no additional cooldown interval is imposed.

Within an episode, each vehicle can receive at most one recourse attempt, and a
vehicle is eligible only if at least two customers remain unserved on its
current route. A trigger therefore denotes \emph{route reconsideration}, not
mandatory rerouting.

The online optimizer cannot change the already executed portion of the route.
At the decision epoch, the following elements are fixed:
the current vehicle node, its executed prefix, its customer assignment, and the
final return to the depot. The decision variable is only the ordering of the
remaining customers assigned to the triggered vehicle.

% ------------------------------------------------------------
\subsection{Conditional local CVaR replanning}
\label{subsec:conditional_replanning}
% ------------------------------------------------------------

Let $r^{\mathrm{inc}}_{k,t_0}$ denote the incumbent residual suffix of triggered
vehicle $k$ at time $t_0$. A conditional scenario bank
$\Omega_{t_0}$ is generated from the observed information described above.
The local online objective uses the same distance and tardiness weights as the
preoperational formulation, with $\rho=\CVaR_{0.95}$.

The incumbent and every candidate suffix are evaluated on the same conditional
scenario bank. Define
\begin{equation}
	G_{t_0}
	=
	\widehat{J}_{\CVaR}
	\left(
	r^{\mathrm{inc}}_{k,t_0};
	\Omega_{t_0}
	\right)
	-
	\widehat{J}_{\CVaR}
	\left(
	r^{\mathrm{sel}}_{k,t_0};
	\Omega_{t_0}
	\right),
	\label{eq:conditional_search_gain}
\end{equation}
where $r^{\mathrm{sel}}_{k,t_0}$ is the best suffix returned by the bounded
search. The suffix is replaced only when the conditional search objective
improves; hence an executed rerouting decision satisfies $G_{t_0}>0$ up to
numerical tolerance.

This is deliberately a \emph{local triggered-vehicle} recourse problem.
It should not be interpreted as exact minimization of the conditional CVaR of
the entire fleet. Final performance, however, is always measured by replaying
the complete fleet execution under the same realized environmental trajectory.

% ------------------------------------------------------------
\subsection{Bounded online search}
\label{subsec:bounded_search}
% ------------------------------------------------------------

Each conditional replanning attempt uses $100$ conditional scenarios.
The ALNS search seeks to evaluate up to $250$ distinct non-null candidate
suffixes but is subject to a hard limit of $25{,}000$ operator decisions.
Null proposals count against the operator-decision budget but require no
scenario evaluation.

If the target number of candidate evaluations is reached first, search terminates
normally. If the operator-decision limit is reached first, the search terminates
at the cap and returns the best admissible suffix found up to that point.
Consequently, if $C\le250$ non-null candidates are actually evaluated, the
conditional scenario-evaluation cost of that attempt is
\begin{equation}
	N_{\mathrm{eval}}
	=
	100(1+C),
	\label{eq:conditional_budget}
\end{equation}
including the incumbent evaluation.

The stopping rule makes online recourse computationally bounded even in residual
states for which the destroy--repair neighbourhood produces very few distinct
feasible alternatives.

% ------------------------------------------------------------
\subsection{Sequential execution algorithm}
\label{subsec:sequential_algorithm}
% ------------------------------------------------------------

Algorithm~\ref{alg:sequential_dddas} summarizes the causal execution protocol.

\begin{algorithm}[t]
	\caption{Causal memory-triggered DDDAS recourse}
	\label{alg:sequential_dddas}
	\begin{algorithmic}[1]
		\State Deploy a preoperational route set $x$
		\State Initialize the Hawkes process and arm the episode trigger
		\While{unserved customers remain}
		\State Select the next vehicle departure in chronological order
		\State Complete any waiting and service at the current node
		\State Set current departure time $t_0$
		\State Observe disruption events with $\tau_q<t_0$
		\State Compute $m(t_0)$ and $\mathrm{OMI}(t_0)$
		\If{$\mathrm{OMI}(t_0)<0.55$}
		\State Rearm the trigger
		\EndIf
		\If{trigger is armed and $\mathrm{OMI}(t_0)\ge0.55$}
		\State Open a high-memory episode and disarm the trigger
		\EndIf
		\If{vehicle is eligible for recourse in the current episode}
		\State Construct a conditional scenario bank from observed information
		\State Run bounded local ALNS--CVaR on the vehicle's residual suffix
		\If{conditional objective improves}
		\State Replace the unexecuted suffix
		\EndIf
		\State Mark the vehicle as attempted for the current episode
		\EndIf
		\State Start the next arc using
		$1+\kappa\,\mathrm{OMI}(t_0)$
		\State Hold that arc multiplier fixed until arrival
		\State Register the realized baseline OD time as observed upon arrival
		\EndWhile
		\State Compute full-fleet realized distance, tardiness, and service outcomes
	\end{algorithmic}
\end{algorithm}

The static and DDDAS executions used for paired evaluation share the same
realized Hawkes path and baseline travel-time realization. Thus differences
between them arise from route decisions rather than from different environmental
draws.

% ------------------------------------------------------------
\subsection{Methodological separation of ex-ante and ex-post quantities}
\label{subsec:exante_expost}
% ------------------------------------------------------------

A final distinction is essential for interpreting online performance.
The quantity in Eq.~\eqref{eq:conditional_search_gain} is an
\emph{ex-ante conditional} comparison based on a distribution of possible
future continuations at the decision epoch. By contrast, realized route cost is
computed after execution from the single future trajectory that actually occurs,
\begin{equation}
	J_{\mathrm{realized}}
	=
	D_{\mathrm{realized}}
	+
	2Z_{\mathrm{realized}},
	\label{eq:realized_objective}
\end{equation}
for the confirmatory sequential experiment.

Accordingly,
\[
G_{t_0}>0
\]
does not imply that the selected suffix must produce a lower
$J_{\mathrm{realized}}$ on every subsequent realization.
The former is a conditional risk statement; the latter is one sample from the
conditional future distribution. Maintaining this separation prevents
decision-time risk improvement from being conflated with guaranteed ex-post
performance.















% ============================================================
\section{Experimental Design and Evaluation}
\label{sec:experiments}
% ============================================================

The experimental design was constructed to separate four questions that are
frequently conflated in dynamic stochastic routing:
(i) whether temporal persistence changes routing risk,
(ii) whether explicit tail-risk optimization improves preoperational routing,
(iii) whether learned operator selection improves upon adaptive search under a
comparable computational budget, and
(iv) whether causal online recourse provides additional value once route
execution has begun.

All stochastic comparisons use common random numbers whenever methods are
evaluated against one another. Search, validation, transfer, and sequential
recourse experiments were frozen before the corresponding confirmatory
comparisons were inspected. Results from earlier exploratory implementations are
not included in the analyses reported below.

% ------------------------------------------------------------
\subsection{Athens-derived operational instances}
\label{subsec:instances}
% ------------------------------------------------------------

The experiments use nine daily last-mile delivery instances derived from the
Athens operational data described by \citet{Vrani2025DataInBrief}.
The source data provide the spatial structure, customer orders, weight and
volume demands, service information, and OD travel-time descriptors used to
construct the routing instances.

The nine days contain between 63 and 84 customer orders, excluding the depot,
and exhibit substantial heterogeneity in daily demand and in the consistency of
the OD travel-time triplets. Table~\ref{tab:instance_characteristics}
summarizes the principal structural characteristics used in the present study.

\begin{table}[t]
	\centering
	\caption{Structural characteristics of the nine Athens-derived daily instances.
		The final column reports the fraction of OD travel-time triplets requiring
		ordering regularization for stochastic PERT sampling. The raw source values
		themselves are not modified.}
	\label{tab:instance_characteristics}
	\begin{tabular}{rrrrr}
		\toprule
		Day &
		Customers &
		Weight (kg) &
		Volume (m$^3$) &
		PERT-reg. arcs (\%) \\
		\midrule
		1 & 78 & 1880.498 & 9.0200 & 7.61 \\
		2 & 63 & 1394.549 & 5.2281 & 23.14 \\
		3 & 67 &  687.633 & 4.0536 & 8.89 \\
		4 & 69 & 1689.368 & 4.7584 & 0.62 \\
		5 & 75 & 1270.167 & 4.4493 & 1.04 \\
		6 & 77 & 1301.434 & 4.5925 & 2.93 \\
		7 & 66 & 1034.561 & 3.6932 & 0.25 \\
		8 & 74 & 1095.370 & 3.2131 & 0.22 \\
		9 & 84 & 1301.537 & 5.3159 & 0.10 \\
		\bottomrule
	\end{tabular}
\end{table}

The empirical routing data and the experimental disruption mechanism play
different roles. The Athens data define the operational base instances, whereas
the tightened service windows and Hawkes disruption process define controlled
stress scenarios. Accordingly, the experiments should be interpreted as
\emph{Athens-derived real operational instances under controlled stress
	scenarios}, not as a reconstruction of historical disruption events in Athens.

The vehicle fleet is fixed at six vehicles. Weight and volume capacities are
experimental design parameters established before the comparative analyses;
they should not be interpreted as observed capacities of a specific Athens
fleet. Capacity feasibility is enforced simultaneously in both dimensions.

For the deterministic OR-Tools computations, continuous demand quantities are
mapped to integer units without truncating positive demands. Weight is scaled by
$10^3$ and volume by $10^6$ before integer optimization. Stochastic evaluation
is subsequently carried out in the original physical units.

The controlled stress protocol uses service-window widths of 180 minutes.
Its purpose is not to reproduce a specific historical operating day exactly,
but to provide a common reliability environment in which differences among
mean-oriented, tail-oriented, learning-augmented, and dynamically adaptive
routing policies can be exposed.

% ------------------------------------------------------------
\subsection{Preoperational route initialization}
\label{subsec:initialization}
% ------------------------------------------------------------

To avoid method-specific advantages arising from different starting solutions,
route initialization is performed independently of the stochastic search
methods.

For the nine-day cross-instance experiment, a common OR-Tools initialization
protocol was frozen before comparative stochastic optimization.
The protocol tests a fixed ordered collection of standard first-solution
strategies, including local cheapest insertion, path-most-constrained arc,
savings, global-cheapest arc, and path-cheapest arc, combined with guided local
search. Each strategy is allocated the same time allowance, and the first
feasible solution generated by the frozen protocol is retained.

This fallback mechanism is applied uniformly across days rather than selecting
the best initializer retrospectively for each instance. In particular, different
strategies may provide the first feasible route on different days without
constituting instance-specific tuning.

The earlier source-day learning validation on Day~1 uses its independently
frozen Day~1 initialization. The subsequent nine-day experiment uses the common
cross-day initialization protocol described above. Numerical results from these
two stages are therefore interpreted within their respective experimental
designs and are not pooled as if they arose from an identical initial route.

% ------------------------------------------------------------
\subsection{Nominal stochastic configuration}
\label{subsec:nominal_configuration}
% ------------------------------------------------------------

The nominal persistent-disruption configuration uses
\begin{equation}
	\beta=0.60,
	\qquad
	\eta=0.35,
	\qquad
	\mu=0.01,
	\qquad
	\kappa=0.90.
	\label{eq:nominal_hawkes_parameters}
\end{equation}
With the parameterization $\xi=\eta\beta$, the nominal excitation amplitude is
$\xi=0.21$, and the stationary mean Hawkes event intensity is
\begin{equation}
	\bar{\lambda}
	=
	\frac{0.01}{1-0.35}
	\simeq
	0.01538.
\end{equation}

The modified PERT model uses
$\lambda_{\mathrm{PERT}}=4$, and all tail-risk experiments use
\begin{equation}
	\alpha=0.95.
\end{equation}
The canonical objective weights are
\begin{equation}
	\lambda_D=1,
	\qquad
	\lambda_T=2,
	\qquad
	\lambda_W=\lambda_V=10.
	\label{eq:objective_weights}
\end{equation}
Because weight and volume feasibility are enforced by the route-construction
mechanism, the last two terms remain zero for reported feasible solutions.

% ------------------------------------------------------------
\subsection{Frozen stochastic evaluation banks}
\label{subsec:frozen_banks}
% ------------------------------------------------------------

Four conceptually distinct scenario-bank roles are used:
\[
\Omega_{\mathrm{train}},
\qquad
\Omega_{\mathrm{val}},
\qquad
\Omega_{\mathrm{search}},
\qquad
\Omega_{\mathrm{test}}.
\]
The corresponding baseline configuration uses
\begin{equation}
	|\Omega_{\mathrm{train}}|
	=
	|\Omega_{\mathrm{val}}|
	=
	|\Omega_{\mathrm{search}}|
	=
	100,
	\qquad
	|\Omega_{\mathrm{test}}|
	=
	1000.
	\label{eq:scenario_bank_sizes}
\end{equation}

Training scenarios are used only when fitting learned policies.
Validation scenarios support policy selection without using the final test bank.
Search scenarios are used to compare candidate and incumbent routes during
optimization. The test bank is held out from search and is used only for final
performance evaluation.

Within any paired comparison, the same test scenarios are supplied to all
methods. Candidate and incumbent solutions within ALNS or learned search also
share the same search bank. Thus stochastic differences between methods are
paired at the scenario level rather than being confounded with independently
generated disruption realizations.

% ------------------------------------------------------------
\subsection{Search-budget matching}
\label{subsec:fair_budget}
% ------------------------------------------------------------

A central objective of the experimental design is to distinguish algorithmic
performance from computational expenditure. For the principal ALNS and frozen
learned-policy comparisons, online search is therefore matched by the number of
candidate-route--scenario evaluations.

Each search uses $100$ search scenarios and evaluates up to $250$ non-null
candidate route sets in addition to the incumbent. The nominal stochastic
evaluation budget is consequently
\begin{equation}
	B_{\mathrm{search}}
	=
	100(1+250)
	=
	25{,}100
	\label{eq:search_budget}
\end{equation}
scenario evaluations.

The same search budget is applied to adaptive ALNS and to the feed-forward and
GRU-guided variants. Runtime is reported as a complementary quantity, but is not
used as the sole fairness criterion because stochastic objective evaluation and
neighbourhood generation have method-dependent computational structure.

Offline RL training is \emph{not} included inside
$B_{\mathrm{search}}$. Instead, it is reported separately as an additional
up-front computational cost. In the source-day experiments, fitting one policy
requires approximately $2.3\times10^5$ scenario evaluations, or roughly an order
of magnitude more stochastic evaluations than one frozen-policy search.
This distinction prevents the cost of policy learning from being hidden inside
a nominally equal online-search comparison.

% ------------------------------------------------------------
\subsection{Source-day learning experiment}
\label{subsec:source_day_rl}
% ------------------------------------------------------------

The first learning experiment is conducted on Day~1 and compares six
risk--search combinations:
\[
\begin{array}{lll}
	\text{ALNS-Mean},      &
	\text{RL-FF-Mean},     &
	\text{RL-GRU-Mean},
	\\
	\text{ALNS-CVaR},      &
	\text{RL-FF-CVaR},     &
	\text{RL-GRU-CVaR}.
\end{array}
\]
Here ``FF'' denotes the feed-forward operator-selection policy and ``GRU''
denotes the recurrent policy in Eq.~\eqref{eq:gru_policy}.

Ten independent paired repetitions are used for final source-day evaluation.
Within each repetition, competing methods share the relevant frozen stochastic
banks. Learned policies are trained and validated before final test evaluation;
the $1000$-scenario final test bank is not used for policy fitting or
model selection.

A uniform operator-selection policy is evaluated as an additional control.
It uses the same search machinery and stochastic budget but samples among the
available destroy operators uniformly rather than through learned or adaptive
selection. This control is important because it distinguishes the value of a
trained policy from improvements that could arise simply through stochastic
operator variation.

The source-day analysis therefore addresses two different questions:
whether learned policies outperform the genuine adaptive ALNS baseline, and
whether any observed learned-policy advantage exceeds that of a simple uniform
operator selector.

% ------------------------------------------------------------
\subsection{Zero-shot cross-day policy transfer}
\label{subsec:transfer}
% ------------------------------------------------------------

Generalization is evaluated without target-day retraining.
Policies are trained only on Day~1 and then transferred directly to Days~2--9.
No target-day learning, fine-tuning, or target-specific adjustment of policy
parameters is permitted.

The transfer experiment uses 40 frozen source-policy checkpoints, corresponding
to four checkpoints from each of ten source training repetitions. Every
checkpoint is evaluated on the eight target days using the same target-day route
initializations and stochastic evaluation protocol as the corresponding
uniform-policy control.

The primary transfer contrasts compare learned and uniform operator selection
for both feed-forward and recurrent policies under the mean- and CVaR-oriented
objectives. The target day, rather than an individual checkpoint or stochastic
run, is the cross-instance unit of analysis. Four prespecified primary transfer
contrasts are corrected jointly using the Holm procedure.

This design tests zero-shot transportability of the learned search policy rather
than its ability to be retrained on each new instance.

% ------------------------------------------------------------
\subsection{Persistence experiment}
\label{subsec:persistence_experiment}
% ------------------------------------------------------------

A separate set of experiments isolates the effect of temporal persistence.
The Hawkes branching ratio is varied over
\begin{equation}
	\eta
	\in
	\{0,\;0.15,\;0.35,\;0.55,\;0.70\}.
	\label{eq:eta_grid}
\end{equation}

A direct change in $\eta$ would simultaneously alter event clustering and the
stationary mean event rate. To prevent this confounding, the first control fixes
the stationary mean intensity at its nominal value,
$\bar{\lambda}$, by setting
\begin{equation}
	\mu_{\eta}
	=
	\bar{\lambda}(1-\eta).
	\label{eq:mu_matched}
\end{equation}
Thus $\eta$ changes the temporal clustering of the disruption process without
changing its stationary mean number of events.

For this persistence experiment, the operational memory transformation that
maps realized events into OMI retains the nominal decay and excitation
parameters. This ensures that the $\eta=0$ generator provides a meaningful
Poisson-like no-self-excitation comparison while the route can still experience
nonzero disruption exposure when external events occur.

Before routing comparisons, $10{,}000$ environmental trajectories are generated
for each persistence regime. These simulations are used only to characterize
the experimental manipulation---including event-count dispersion, no-event
frequency, event timing, and time-averaged OMI---rather than to tune routing
algorithms.

% ------------------------------------------------------------
\subsection{Mean-impact-matched persistence control}
\label{subsec:mean_impact_control}
% ------------------------------------------------------------

Matching event intensity does not by itself guarantee equal mean operational
impact because changes in clustering can alter the time profile of OMI.
We therefore introduce a second, stricter control that matches the
\emph{continuous-time mean environmental travel-time multiplier}
\begin{equation}
	\mathbb{E}_t
	\left[
	1+\kappa_{\eta}\mathrm{OMI}_{\eta}(t)
	\right].
	\label{eq:mean_multiplier_control}
\end{equation}

The nominal target is
\[
\mathbb{E}_t[M(t)]
=
1.052489.
\]
The resulting prospectively computed impact coefficients are
\begin{equation}
	\begin{array}{c|ccccc}
		\eta
		&0
		&0.15
		&0.35
		&0.55
		&0.70\\
		\hline
		\kappa_{\eta}
		&0.6953
		&0.7646
		&0.9000
		&1.1418
		&1.4816.
	\end{array}
	\label{eq:kappa_matched_values}
\end{equation}

This control matches the continuous-time environmental mean multiplier, not the
mean multiplier sampled by any particular route. Route-sampled exposure remains
route-dependent because methods depart on arcs at different times.

The confirmatory persistence-routing experiment is performed on Day~1 using ten
paired repetitions per persistence regime. Candidate search uses $100$
scenarios and final evaluation uses $1000$ scenarios. No search parameter is
retuned as $\eta$ changes.

The experiment evaluates both the direct change in routing tail risk as
persistence increases and whether persistence systematically changes the
relative advantage of CVaR-oriented over mean-oriented ALNS.

% ------------------------------------------------------------
\subsection{Nine-day risk-aware routing experiment}
\label{subsec:nine_day_experiment}
% ------------------------------------------------------------

Cross-instance robustness of explicit tail-risk optimization is assessed on all
nine Athens-derived days under the nominal disruption configuration.
Three paired optimization repetitions are performed on each day using the
common cross-day initialization protocol.

The principal comparison is
\[
\text{ALNS-Mean}
\quad\text{versus}\quad
\text{ALNS-CVaR}.
\]
Both methods use the same adaptive ALNS engine, neighbourhoods, search length,
scenario-bank size, and final test bank. The comparison therefore isolates the
effect of changing the optimized risk functional from expected tardiness to
$\CVaR_{0.95}$.

For each repetition, candidate search is based on $100$ stochastic scenarios and
final evaluation on $1000$ held-out scenarios. The three optimization
replications are first averaged within each operational day. Consequently, the
cross-instance inferential unit is the \emph{day}, giving $n=9$ independent
day-level contrasts rather than treating the 27 optimization runs as independent
instances.

The two prespecified primary contrasts are
\begin{align}
	\Delta J_{\mathrm{CVaR}}
	&=
	J_{\mathrm{CVaR}}(\text{ALNS-Mean})
	-
	J_{\mathrm{CVaR}}(\text{ALNS-CVaR}),
	\label{eq:delta_j_cvar}
	\\
	\Delta T_{\mathrm{tail}}
	&=
	\CVaR_{0.95}(Z_{\text{ALNS-Mean}})
	-
	\CVaR_{0.95}(Z_{\text{ALNS-CVaR}}).
	\label{eq:delta_tail}
\end{align}
Positive values favor CVaR-oriented search.

Across-day uncertainty is summarized using bootstrap confidence intervals over
the nine day-level contrasts. Because the direction of the paired contrast is
prespecified, an exact two-sided sign-flip test is additionally used as a
distribution-free check of cross-day consistency.

% ------------------------------------------------------------
\subsection{Confirmatory sequential DDDAS experiment}
\label{subsec:r5h_design}
% ------------------------------------------------------------

The online DDDAS experiment is deliberately separated from the nine-day
preoperational analysis. Its purpose is not to establish cross-instance
superiority over the entire dataset, but to examine the mechanics and
realized consequences of causal sequential recourse in greater depth.

Three representative instances---Days~2, 5, and 7---were selected
prospectively from the nine-day set using structural information only, before
the confirmatory sequential outcomes were inspected. Day~1 was excluded from
this selection because it had already been used extensively during development
and runtime validation.

For each selected day, the three independently generated ALNS-CVaR routes from
the nine-day experiment are frozen as preoperational starting solutions.
Twenty realized environmental trajectories are then generated per day.
The same realized Hawkes and baseline travel-time trajectory is replayed across
the three preoperational route replicates for a given day--path pair, producing
\[
3\ \text{days}
\times
3\ \text{route replicates}
\times
20\ \text{realized paths}
=
180
\]
paired static-versus-DDDAS executions.

Because the same realized environment is shared across the three route
replicates, these 180 runs do not constitute 180 independent environmental
observations. There are instead
\[
3\times20=60
\]
unique day--realized-path environments.
Recourse-search random seeds are distinct across route replicates so that the
online optimizer itself does not reuse an identical random search trajectory.

Each DDDAS recourse attempt uses $100$ conditional scenarios and the bounded
stopping rule in Section~\ref{subsec:bounded_search}.
Static and DDDAS execution share the same complete realized environmental path.
The static comparator simply continues executing its frozen preoperational route,
whereas DDDAS may alter eligible residual suffixes using only causally observed
information.

The realized composite criterion is
\begin{equation}
	J_{\mathrm{realized}}
	=
	D_{\mathrm{realized}}
	+
	2Z_{\mathrm{realized}},
\end{equation}
and paired contrasts are defined as
\begin{align}
	\Delta J
	&=
	J_{\mathrm{static}}
	-
	J_{\mathrm{DDDAS}},
	\\
	\Delta Z
	&=
	Z_{\mathrm{static}}
	-
	Z_{\mathrm{DDDAS}},
	\\
	\Delta D
	&=
	D_{\mathrm{static}}
	-
	D_{\mathrm{DDDAS}}.
\end{align}
Positive values favor DDDAS.

For scientific summaries, the three route replicates are first averaged within
each fixed day--realized-path pair. This yields 20 Monte Carlo path-level
contrasts per selected day. Within-day uncertainty is summarized by
nonparametric bootstrap confidence intervals over these 20 realized
environments, conditional on the frozen three-route ensemble.

Only three days are examined in this deep sequential experiment. We therefore
do not use the 180 runs, or even the 60 path-level observations, as independent
cross-instance replicates. Cross-day R5-H summaries are descriptive rather than
population-level hypothesis tests. Cross-instance evidence is instead provided
by the nine-day preoperational experiment in
Section~\ref{subsec:nine_day_experiment}.

No ``realized CVaR'' is constructed from an individual executed trajectory.
The decision-time CVaR used by the online optimizer is a conditional
distributional quantity, whereas each static--DDDAS comparison produces one
subsequently realized outcome.

% ------------------------------------------------------------
\subsection{Performance measures}
\label{subsec:performance_metrics}
% ------------------------------------------------------------

The principal routing measures are:
\begin{enumerate}
	\item total route distance $D$;
	\item mean realized total tardiness $\mathbb{E}[Z]$;
	\item $\CVaR_{0.95}(Z)$;
	\item fraction of customers served on time;
	\item maximum realized customer tardiness; and
	\item the objective corresponding to the risk functional used in the
	relevant experiment.
\end{enumerate}

Route-dependent mean and maximum OMI exposure are recorded as descriptive
diagnostics. These quantities summarize the disruption states sampled at the
actual departure times of the route and must not be interpreted as changes in
the environmental Hawkes process itself.

For learning experiments, computational accounting includes online
candidate--scenario evaluations, search runtime, and offline training cost.
For sequential DDDAS, additional operational measures include the number of
high-memory episodes, recourse attempts, executed route changes, evaluated
non-null candidates, operator decisions, decision-cap terminations, conditional
scenario evaluations, and recourse runtime.

A DDDAS trigger is therefore evaluated both as a computational event and as an
operational decision. The opening of a high-memory episode does not imply that
a route change will occur.

% ------------------------------------------------------------
\subsection{Paired analysis and inferential units}
\label{subsec:inference}
% ------------------------------------------------------------

The experimental design consistently distinguishes stochastic repetitions from
independent operational instances.

For the Day~1 learning study, ten paired stochastic repetitions quantify
source-instance variability.
For zero-shot transfer, the target operational day is the cross-instance unit
($n=8$).
For the principal ALNS-Mean versus ALNS-CVaR comparison, the inferential unit is
again the operational day ($n=9$), after averaging the three optimization
replications within each day.
For the deep DDDAS study, uncertainty is summarized separately within each of
the three prospectively selected days, using the 20 realized environmental
paths after averaging the three frozen route replicates.

This hierarchy avoids treating algorithmic seeds or common-random-number
replications as independent real-world instances.

All reported method differences retain a fixed sign convention.
For the nine-day risk experiment,
\[
\text{ALNS-Mean}-\text{ALNS-CVaR}>0
\]
favors CVaR-oriented optimization.
For DDDAS,
\[
\text{static}-\text{DDDAS}>0
\]
favors dynamic recourse.
For learned-versus-uniform transfer comparisons, the contrast is defined as
\[
\text{learned}-\text{uniform},
\]
so negative objective differences favor the learned policy.

% ------------------------------------------------------------
\subsection{Computational environment and reproducibility}
\label{subsec:computational_environment}
% ------------------------------------------------------------

All final computations were executed on a single Apple-silicon workstation with
an Apple M5 Max processor, 18 CPU cores, and 64~GB of memory. The computational
pipeline uses Python~3.12.13, NumPy~1.26.4, pandas~3.0.5, SciPy~1.13.1,
PyTorch~2.2.2, and Google OR-Tools~9.10.4067.

The recurrent models are sufficiently small that GPU acceleration does not
provide an advantage on the test platform; final learning and routing
experiments therefore use CPU execution. Runtime comparisons are consequently
made within a common hardware and software environment.

Scenario seeds, train/validation/search/test separation, route initializations,
method configurations, and confirmatory sequential schedules are stored as
reproducibility metadata. Common-random-number checks and sequential replay
tests verify that static and dynamic methods observe identical realized
environmental trajectories whenever a paired comparison requires them.










% =====================
%        RESULTS
% =====================



% ============================================================
\section{Results}
\label{sec:results}
% ============================================================

The results are organized around the four questions posed in
Section~\ref{sec:experiments}. We first verify that the persistence manipulation
changes the temporal organization of disruption events rather than merely their
mean frequency. We then examine how this persistence affects routing risk under
the mean-impact-matched control. Next, we assess the cross-instance value of
explicit CVaR optimization and the incremental contribution of learned operator
selection. Finally, we evaluate the causal sequential DDDAS mechanism on the
three prospectively selected representative days.

% ------------------------------------------------------------
\subsection{Persistence reorganizes disruption exposure into intermittent bursts}
\label{subsec:results_environment}
% ------------------------------------------------------------

The environmental simulation confirms that increasing the Hawkes branching
ratio changes the temporal organization of disruptions substantially even when
the stationary mean event intensity is held fixed through
Eq.~\eqref{eq:mu_matched}.
Across $10{,}000$ simulated environmental trajectories per regime, the mean
number of events remained approximately constant, whereas event-count
dispersion increased sharply with persistence.

Table~\ref{tab:environmental_persistence} summarizes the principal diagnostics.
The event-count Fano factor increased from approximately $1.02$ at $\eta=0$ to
$11.02$ at $\eta=0.70$. At the same time, the probability of observing no
event over the operating horizon increased from $0.03\%$ to $9.17\%$.
The probability of an event occurring within the first 60 minutes decreased
from $0.608$ to $0.238$.

\begin{table}[t]
	\centering
	\caption{Environmental diagnostics for the stationary-rate-matched persistence
		experiment. Each row is based on $10{,}000$ simulated Hawkes trajectories.
		Increasing $\eta$ changes temporal clustering while the stationary mean event
		rate is held fixed.}
	\label{tab:environmental_persistence}
	\begin{tabular}{rrrrr}
		\toprule
		$\eta$ &
		Fano factor &
		$P(N=0)$ &
		$P(\tau_1\le60)$ &
		Mean time-OMI \\
		\midrule
		0.00 & 1.02  & 0.0003 & 0.6081 & 0.07550 \\
		0.15 & 1.40  & 0.0014 & 0.5432 & 0.06865 \\
		0.35 & 2.44  & 0.0052 & 0.4528 & 0.05832 \\
		0.55 & 4.77  & 0.0282 & 0.3437 & 0.04597 \\
		0.70 & 11.02 & 0.0917 & 0.2376 & 0.03543 \\
		\bottomrule
	\end{tabular}
\end{table}

The declining time-averaged OMI should not be interpreted as evidence that a
more persistent environment is ``safer.'' Rather, higher persistence
concentrates disruption activity into more strongly clustered episodes separated
by longer quiet periods. Consistently with this interpretation, the upper tail
of instantaneous OMI remains high across regimes, with its $0.99$ quantile close
to $0.94$ throughout the persistence grid.

These environmental diagnostics are summarized in Fig.~\ref{fig:persistence_environment}.
Thus, the persistence manipulation primarily changes the \emph{temporal
	structure} of environmental exposure: diffuse disruption activity at low
$\eta$ is replaced by increasingly intermittent and clustered bursts as
$\eta$ increases.

% ------------------------------------------------------------
\begin{figure*}[t]
\centering
\includegraphics[width=\textwidth]{figures/main/pdf/Fig2_environment_v3.pdf}
\caption{\textbf{Persistence reorganizes disruption activity under matched stationary event intensity.} (a) Event-count overdispersion rises sharply with the Hawkes branching ratio; the dashed line marks the Poisson reference (Fano factor $=1$). (b) Higher persistence produces more event-free horizons and fewer early events, consistent with longer quiet intervals separated by clustered disruption episodes. (c) Time-averaged OMI decreases as exposure becomes more intermittent; this lower mean should not be interpreted as lower upper-tail routing risk. Environmental diagnostics use 10,000 trajectories per persistence regime.}
\label{fig:persistence_environment}
\end{figure*}

\subsection{Persistence increases tail risk under matched mean environmental impact}
\label{subsec:results_persistence}
% ------------------------------------------------------------

The mean-impact-matched experiment isolates the operational consequences of
this change in temporal structure.
The impact coefficient $\kappa_{\eta}$ is adjusted according to
Eq.~\eqref{eq:kappa_matched_values} so that the continuous-time mean
environmental travel-time multiplier remains fixed across persistence regimes.

Despite this control, routing tail risk increases markedly with persistence.
For the deterministic OR-Tools route, mean tardiness increases from
$13.42$ at $\eta=0$ to $17.19$ at $\eta=0.70$, while
$\CVaR_{0.95}$ of total tardiness increases from $60.74$ to $113.24$.
Across the ten paired repetitions, the increase in OR-Tools tail tardiness
between $\eta=0$ and $\eta=0.70$ is
\begin{equation}
	\Delta\CVaR_{0.95}
	=
	52.503,
	\qquad
	95\%~\mathrm{CI}
	=
	[49.140,\;55.932],
\end{equation}
with all ten paired differences positive
(exact two-sided sign-flip $p=0.00195$).
Mean tardiness increases by $3.771$ minutes
($95\%$ CI $[3.408,4.125]$, $p=0.00195$).

A detectable increase already occurs at the nominal persistence level:
relative to $\eta=0$, OR-Tools $\CVaR_{0.95}$ increases by $12.777$ at
$\eta=0.35$, again with all ten paired contrasts positive
($p=0.00195$).

Figure~\ref{fig:tailrisk_persistence} visualizes the resulting upper-tail increase.
Table~\ref{tab:persistence_routing} shows the corresponding aggregate
performance of the three preoperational routing strategies at selected
persistence levels.

\begin{table*}[t]
	\centering
	\caption{Routing performance under the mean-environmental-impact-matched
		persistence control. Values are averages over ten paired repetitions.
		$J_{\mathrm{CVaR}}=D+2\CVaR_{0.95}(Z)$.}
	\label{tab:persistence_routing}
		\footnotesize
	\begin{tabular}{llrrrr}
		\toprule
		$\eta$ &
		Method &
		Distance &
		Mean tardiness &
		$\CVaR_{0.95}$ tardiness &
		$J_{\mathrm{CVaR}}$ \\
		\midrule
		0.00 &
		OR-Tools &
		367.20 & 13.415 & 60.741 & 488.682 \\
		&
		ALNS-Mean &
		371.76 & 1.550 & 15.780 & 403.320 \\
		&
		ALNS-CVaR &
		376.62 & 0.599 & 8.228 & 393.075 \\
		\midrule
		0.35 &
		OR-Tools &
		367.20 & 14.249 & 73.518 & 514.236 \\
		&
		ALNS-Mean &
		373.28 & 1.468 & 16.580 & 406.441 \\
		&
		ALNS-CVaR &
		378.22 & 0.801 & 8.701 & 395.622 \\
		\midrule
		0.70 &
		OR-Tools &
		367.20 & 17.187 & 113.244 & 593.688 \\
		&
		ALNS-Mean &
		373.95 & 2.924 & 33.443 & 440.836 \\
		&
		ALNS-CVaR &
		378.21 & 1.872 & 21.612 & 421.434 \\
		\bottomrule
	\end{tabular}
\end{table*}

The increase in tail risk cannot be attributed simply to a larger average
environmental multiplier. For the unchanged OR-Tools route, for example, the
route-sampled mean travel-time multiplier is approximately $1.0483$ at
$\eta=0$ and $1.0458$ at $\eta=0.70$, despite the sharp increase in
$\CVaR_{0.95}$. This pattern implicates the temporal concentration of
disruption exposure rather than a larger average route-level multiplier.

The relative advantage of CVaR-oriented search does not, however, increase
robustly with persistence. Define
\[
G(\eta)
=
J_{\mathrm{CVaR}}(\text{ALNS-Mean})
-
J_{\mathrm{CVaR}}(\text{ALNS-CVaR}).
\]
Although the point estimate increases between $\eta=0$ and $\eta=0.70$,
the paired interaction contrast is
\begin{equation}
	G(0.70)-G(0)
	=
	9.157,
	\qquad
	95\%~\mathrm{CI}
	=
	[-15.230,\;33.310],
\end{equation}
with six positive and four negative paired differences
($p=0.494$).
Persistence therefore produces a strong direct increase in tail risk, but the
present data do not support the stronger claim that persistence systematically
amplifies the \emph{relative} benefit of CVaR-oriented search.

% ------------------------------------------------------------
\begin{figure}[t]
\centering
\includegraphics[width=0.667\textwidth]{figures/main/pdf/Fig3_tailrisk_persistence_v3.pdf}
\caption{\textbf{Routing tail risk under the mean-environmental-impact-matched persistence control.} Upper-tail tardiness rises strongly with persistence even though the continuous-time mean environmental travel-time multiplier is matched across regimes. OR-Tools uses the same deterministic route across persistence levels; ALNS-Mean and ALNS-CVaR use the same adaptive search engine but optimize different risk functionals. Values are averages over ten paired repetitions.}
\label{fig:tailrisk_persistence}
\end{figure}

\subsection{CVaR-oriented ALNS improves tail-risk performance across all nine days}
\label{subsec:results_nineday}
% ------------------------------------------------------------

The principal cross-instance comparison evaluates ALNS-Mean and ALNS-CVaR
under the nominal disruption configuration on all nine Athens-derived
operational days.
The two methods use the same adaptive search engine, initial routes,
neighbourhoods, stochastic search budget, and held-out test scenarios.
Only the risk functional used during optimization differs.

The result is highly consistent across instances.
Table~\ref{tab:nineday_gains} reports the day-level paired gains, after averaging
the three optimization repetitions within each day.
Positive values favor ALNS-CVaR.

\begin{table}[t]
	\centering
	\caption{Day-level gains of ALNS-CVaR relative to ALNS-Mean under the nominal
		persistent-disruption configuration. Positive values favor CVaR-oriented
		optimization.}
	\label{tab:nineday_gains}
	\begin{tabular}{rrr}
		\toprule
		Day &
		$\Delta J_{\mathrm{CVaR}}$ &
		$\Delta\CVaR_{0.95}(Z)$ \\
		\midrule
		1 & 10.545 &  6.789 \\
		2 &  4.960 &  4.097 \\
		3 &  1.832 &  1.483 \\
		4 & 11.609 & 12.555 \\
		5 &  1.773 &  5.570 \\
		6 &  8.968 &  4.801 \\
		7 & 10.856 &  9.428 \\
		8 & 13.562 & 11.531 \\
		9 & 13.396 &  9.231 \\
		\bottomrule
	\end{tabular}
\end{table}

For the tail-oriented objective,
\begin{equation}
	\overline{\Delta J_{\mathrm{CVaR}}}
	=
	8.611,
\end{equation}
with day-level standard deviation $4.628$, median $10.545$, and bootstrap
$95\%$ confidence interval
\begin{equation}
	[5.662,\;11.314].
\end{equation}
All nine day-level contrasts are positive, giving an exact two-sided sign-flip
$p=0.00391$.

The same cross-instance consistency is observed directly for tail tardiness:
\begin{equation}
	\overline{\Delta\CVaR_{0.95}(Z)}
	=
	7.276,
\end{equation}
with standard deviation $3.663$, median $6.789$, and bootstrap $95\%$
confidence interval
\begin{equation}
	[5.027,\;9.522].
\end{equation}
Again, all nine day-level contrasts favor ALNS-CVaR
($p=0.00391$).

Aggregating descriptively across the nine days, ALNS-CVaR attains mean distance
$406.55$, mean tardiness $1.92$, tail tardiness $15.84$, and on-time service
fraction $0.9945$. The corresponding ALNS-Mean values are $400.61$, $3.24$,
$23.12$, and $0.9911$, respectively. Thus CVaR optimization accepts a modest
increase in route distance---approximately $5.94$ distance units on average---in
exchange for lower mean and substantially lower upper-tail tardiness.

This trade-off is reflected in the aligned tail-risk objective:
the average $J_{\mathrm{CVaR}}$ is $438.23$ for ALNS-CVaR compared with
$446.85$ for ALNS-Mean. The deterministic OR-Tools route is shorter on average
($398.87$), but produces substantially greater mean tardiness ($11.28$) and
tail tardiness ($47.39$), yielding average
$J_{\mathrm{CVaR}}=493.65$.

Figure~\ref{fig:crossday_cvar} shows the day-level consistency of both prespecified contrasts.
The nine-day result is therefore not that CVaR optimization dominates every
individual performance component. Rather, it consistently trades a small
distance increase for a sufficiently large reduction in upper-tail service
failure to improve the prespecified risk-oriented objective on every operational
day.

% ------------------------------------------------------------
\begin{figure*}[t]
\centering
\includegraphics[width=\textwidth]{figures/main/pdf/Fig4_crossday_v3.pdf}
\caption{\textbf{Cross-instance gains from explicit CVaR-oriented planning.} Day-level contrasts compare ALNS-Mean with ALNS-CVaR after averaging the three paired optimization repetitions within each operational day; positive values favour CVaR-oriented search. (a) The aligned tail-risk objective improves on all nine days. (b) $\mathrm{CVaR}_{0.95}$ tardiness likewise improves on all nine days. The diamond and dashed vertical line report the across-day mean; the horizontal bar is the bootstrap 95\% confidence interval.}
\label{fig:crossday_cvar}
\end{figure*}

\subsection{Learned operator selection does not provide a robust incremental advantage}
\label{subsec:results_rl}
% ------------------------------------------------------------

The Day~1 learning experiment evaluates whether feed-forward or recurrent
operator selection improves on the genuine adaptive ALNS backbone when the
online stochastic search budget is held fixed.

Table~\ref{tab:rl_day1} reports the ten-repetition source-day results.
For mean-oriented optimization, the three methods produce closely overlapping
performance. For CVaR-oriented optimization, the learned policies have somewhat
lower point estimates of the aligned objective than ALNS-CVaR, but the
between-method uncertainty remains substantial.

\begin{table*}[t]
	\centering
	\caption{Day~1 source-instance learning results over ten paired repetitions.
		Values are mean $\pm$ standard deviation. The final column reports the objective
		aligned with the criterion being optimized: $J_{\mathrm{mean}}$ for
		mean-oriented rows and $J_{\mathrm{CVaR}}$ for CVaR-oriented rows.}
	\label{tab:rl_day1}
		\footnotesize
	\begin{tabular}{llrrr}
		\toprule
		Criterion &
		Method &
		Mean tardiness &
		$\CVaR_{0.95}$ tardiness &
		Aligned objective \\
		\midrule
		Mean &
		ALNS &
		$1.440\pm1.017$ &
		$17.407\pm10.885$ &
		$374.970\pm3.240$ \\
		&
		RL-FF &
		$1.783\pm1.489$ &
		$19.082\pm14.064$ &
		$375.585\pm3.629$ \\
		&
		RL-GRU &
		$1.355\pm0.951$ &
		$16.322\pm10.320$ &
		$375.479\pm3.225$ \\
		\midrule
		CVaR &
		ALNS &
		$0.812\pm0.685$ &
		$9.413\pm5.879$ &
		$396.437\pm13.213$ \\
		&
		RL-FF &
		$0.515\pm0.334$ &
		$7.785\pm3.995$ &
		$393.670\pm12.277$ \\
		&
		RL-GRU &
		$0.708\pm0.741$ &
		$8.501\pm7.323$ &
		$393.181\pm15.841$ \\
		\bottomrule
	\end{tabular}
\end{table*}

Figure~\ref{fig:rl_validation} places these source-day point estimates beside the separate offline training burden.
The paired contrasts do not establish a robust advantage of either learned
policy over adaptive ALNS: the corresponding confidence intervals include zero.
Likewise, learned operator selection does not significantly outperform the
uniform operator-selection control under the matched search budget; the
associated source-day comparisons all have $p>0.59$.

The lack of a source-day significance result is reinforced by the zero-shot
transfer experiment.
Policies trained only on Day~1 are applied without fine-tuning to Days~2--9.
Across the four prespecified learned-versus-uniform contrasts---feed-forward and
GRU policies under mean- and CVaR-oriented objectives---none remains significant
after Holm correction at the target-day level.

These results do not imply that learning-based operator selection is generally
ineffective. They show instead that, under the present state representation,
training protocol, and matched online-search budget, the learned selectors do
not demonstrate a reproducible incremental advantage over either genuine
adaptive ALNS or simple uniform operator selection.

The recurrent policy also fails to exhibit a consistent advantage over the
feed-forward alternative. Thus, although the environment contains temporal
persistence, recurrence itself does not guarantee improved search control once
the relevant environmental and route-exposure summaries are already represented
in the observable state.

This conclusion is especially relevant when computational cost is considered.
Training one policy requires approximately $2.3\times10^5$ scenario
evaluations, around $9$--$10$ times the stochastic evaluation cost of a single
frozen-policy search. The offline learning cost therefore represents a
substantial additional expenditure that is not justified by a corresponding
robust performance gain in the present experiments.

% ------------------------------------------------------------
\begin{figure*}[t]
\centering
\includegraphics[width=\textwidth]{figures/main/pdf/Fig5_RL_v3.pdf}
\caption{\textbf{Source-day performance and computational burden of learned operator selection.} (a,b) Points show the mean aligned objective and error bars one standard deviation over ten paired Day~1 repetitions; the dotted line marks the ALNS mean. These panels summarize descriptive source-day variability; the paired inferential comparisons reported in the text do not establish a robust learned-policy advantage. (c) Offline policy fitting requires approximately $2.3\times10^5$ scenario evaluations per policy, about 9.2 times the stochastic evaluation budget of one frozen-policy search.}
\label{fig:rl_validation}
\end{figure*}

\subsection{Sequential DDDAS acts as selective rather than systematic recourse}
\label{subsec:results_dddas}
% ------------------------------------------------------------

The confirmatory DDDAS experiment evaluates sequential recourse on Days~2, 5,
and 7 using $180$ paired static-versus-DDDAS executions.
These correspond to $60$ unique day--environment combinations, with three
frozen preoperational ALNS-CVaR route replicates evaluated against each
realized environmental trajectory.

Table~\ref{tab:dddas_scientific} reports the principal paired outcomes after
first averaging the three route replicates within each day--realized-path unit.
Positive differences favor DDDAS.

\begin{table*}[t]
	\centering
	\caption{Sequential DDDAS realized-performance contrasts. Values are means over
		20 realized environmental paths after averaging the three frozen
		preoperational-route replicates within each path. Parentheses give
		path-bootstrap $95\%$ confidence intervals. Positive values favor DDDAS.}
	\label{tab:dddas_scientific}
		\footnotesize
	\begin{tabular}{rrrr}
		\toprule
		Day &
		$\Delta J_{\mathrm{realized}}$ &
		$\Delta Z_{\mathrm{realized}}$ &
		$\Delta D_{\mathrm{realized}}$ \\
		\midrule
		2 &
		$-0.169\;[-0.472,\;0.182]$ &
		$ 0.058\;[0.000,\;0.173]$ &
		$-0.285\;[-0.505,\;-0.087]$ \\
		5 &
		$ 0.087\;[-0.377,\;0.549]$ &
		$ 0.056\;[-0.131,\;0.257]$ &
		$-0.025\;[-0.300,\;0.202]$ \\
		7 &
		$ 0.122\;[-1.481,\;2.200]$ &
		$ 0.416\;[-0.298,\;1.396]$ &
		$-0.710\;[-1.100,\;-0.372]$ \\
		\bottomrule
	\end{tabular}
\end{table*}

There is no systematic improvement in the realized composite objective.
The mean of the three day-level $\Delta J$ values is only $0.013$, with
Day~2 favoring static execution and Days~5 and 7 showing small positive mean
differences whose bootstrap intervals include zero.

Distance displays a clearer pattern.
DDDAS increases realized distance on average on all three days:
\[
\Delta D
=
-0.285,\;-0.025,\;-0.710
\]
for Days~2, 5, and 7, respectively.
For Days~2 and 7, the entire path-bootstrap interval is below zero.
Mean tardiness moves in the opposite direction on all three days, but the
estimated improvements are small relative to their uncertainty.

Day~7 illustrates the distinction particularly clearly.
Its mean realized objective difference is slightly positive
($0.122$), whereas the median is negative ($-0.033$).
A leave-one-realized-path-out sensitivity analysis gives
$\Delta J$ means ranging from $-0.717$ to $0.575$.
The positive mean is therefore strongly influenced by a rare path on which
DDDAS avoids substantial tardiness. No observation is removed from the reported
analysis; the leave-one-out calculation is used only to characterize
sensitivity of the descriptive mean.

% ------------------------------------------------------------
\subsection{Most DDDAS triggers leave the incumbent route unchanged}
\label{subsec:results_dddas_operations}
% ------------------------------------------------------------

The online mechanism is markedly conservative in its actual routing actions.
Across the 180 sequential executions, the controller performs 523 recourse
attempts but only 87 result in a route modification:
\begin{equation}
	\frac{87}{523}
	=
	16.6\%.
\end{equation}
The remaining attempts retain the incumbent residual route.

At the run level, 115 of the 180 static--DDDAS comparisons contain no route
change at all. For these 115 runs, static and DDDAS distance, tardiness, and
realized objective are identical to numerical precision.
Only 65 runs contain one or more actual route changes.

Table~\ref{tab:dddas_operational} summarizes the operational behavior by day.

\begin{table}[t]
	\centering
	\caption{Operational behavior of the bounded DDDAS controller.}
	\label{tab:dddas_operational}
	\begin{tabular}{rrrrrr}
		\toprule
		Day &
		Runs &
		Episodes &
		Attempts &
		Changes &
		Cap terminations \\
		\midrule
		2 & 60 & 110 & 151 & 24 & 62 \\
		5 & 60 & 151 & 199 & 23 & 83 \\
		7 & 60 & 130 & 173 & 40 & 49 \\
		\midrule
		Total &
		180 & 391 & 523 & 87 & 194 \\
		\bottomrule
	\end{tabular}
\end{table}

The $25{,}000$-operator-decision cap is reached in 194 attempts,
or $37.1\%$ of all recourse attempts.
Only three of these capped searches ultimately produce a route change.
Thus, in the great majority of cap-limited residual states, the bounded search
exhausts its operator-decision budget without identifying a suffix that
improves sufficiently on the incumbent conditional objective.

Across all 523 recourse attempts, the sequential experiment performs
$8{,}393{,}300$ conditional scenario evaluations.
Total measured DDDAS execution time over the 180 runs is $0.387$ hours on the
reported computational platform, corresponding to approximately $7.74$ seconds
per complete run on average. These timings reflect the optimized sequential
implementation used for the confirmatory experiment.

The controller also frequently modifies a later part of the route rather than
the immediately following arc. Only 19 of the 87 executed route modifications
change the next customer visited at the decision epoch. The remaining changes
preserve the immediate next node and modify a later portion of the residual
suffix.

These observations reinforce the operational interpretation of the OMI trigger:
crossing the threshold initiates \emph{reconsideration} of the current route,
not compulsory rerouting.

% ------------------------------------------------------------
\subsection{Conditional risk improvement does not imply realized-cost improvement}
\label{subsec:results_exante_expost}
% ------------------------------------------------------------

The sequential logs permit direct comparison between the objective used when a
rerouting decision is made and the subsequently realized outcome.

All 87 executed route modifications satisfy
\begin{equation}
	G_{t_0}
	=
	\widehat{J}_{\CVaR}
	(r^{\mathrm{inc}}_{t_0})
	-
	\widehat{J}_{\CVaR}
	(r^{\mathrm{sel}}_{t_0})
	>
	0,
\end{equation}
so every executed change improves the local conditional CVaR objective on the
decision-time scenario bank.

This consistency does not translate into guaranteed ex-post improvement.
Among the 65 complete executions containing at least one route modification,
the final realized composite objective favors DDDAS in 18 runs, is equal within
a numerical tolerance of $10^{-12}$ in 3 runs, and favors static execution in
44 runs.

Therefore,
\begin{equation}
	G_{t_0}>0
	\quad\not\Rightarrow\quad
	J_{\mathrm{DDDAS,realized}}
	<
	J_{\mathrm{static,realized}}.
\end{equation}

This is not a contradiction between the optimization and evaluation criteria.
The first quantity compares two residual routes over a conditional distribution
of possible future trajectories at the decision epoch. The second compares the
two policies after one particular future trajectory has been realized.
A route that reduces conditional tail risk can therefore incur a distance
premium without receiving the corresponding insurance benefit on a specific
realization.

Consistently with this interpretation, there is no stable descriptive
relationship across the three days between the magnitude of the local
conditional search gain and the final realized objective difference.
The sequential results are therefore better characterized as evidence of
\emph{selective risk-sensitive recourse} than as evidence of systematic
realized-cost improvement.

% ------------------------------------------------------------
\begin{figure*}[t]
\centering
\includegraphics[width=\textwidth]{figures/main/pdf/Fig6_DDDAS_v3.pdf}
\caption{\textbf{Realized effects and operational selectivity of causal DDDAS recourse.} Static-minus-DDDAS contrasts are shown for (a) realized composite objective, (b) realized tardiness and (c) realized distance; positive values favour DDDAS. Points are means over 20 realized environmental paths after averaging the three frozen route replicates; error bars are path-bootstrap 95\% confidence intervals. (d) Attempt-level and run-level outcomes are shown as proportions of their own denominators. Only 87 of 523 recourse attempts modify a route, and 115 of 180 complete runs contain no route change. Every executed route modification improves the local conditional search objective, but realized full-run cost does not improve systematically: among the 65 changed runs, 18 improve, 3 tie within $10^{-12}$, and 44 worsen relative to static execution.}
\label{fig:dddas}
\end{figure*}

\subsection{Integrated interpretation}
\label{subsec:results_integrated}
% ------------------------------------------------------------

Figure~\ref{fig:dddas} summarizes the realized DDDAS contrasts and its strong operational selectivity.
Taken together, the experiments separate three effects that would otherwise be
easy to conflate.

First, temporal persistence has an independent operational consequence:
even after controlling the continuous-time mean environmental travel-time
multiplier, greater clustering substantially increases upper-tail tardiness.

Second, explicitly accounting for this tail risk at the preoperational planning
stage produces the most robust algorithmic improvement observed in the study.
ALNS-CVaR improves both the aligned tail-risk objective and
$\CVaR_{0.95}$ tardiness relative to ALNS-Mean on all nine operational days.

Third, additional algorithmic complexity does not automatically yield further
benefit. Learned operator selection does not robustly outperform adaptive or
uniform operator selection under the matched evaluation protocol, and causal
DDDAS recourse does not systematically reduce realized composite cost despite
making locally risk-improving decisions whenever a route modification is
executed.

The resulting empirical picture is therefore not one of uniform algorithmic
dominance. Rather, the strongest evidence supports explicit
\emph{risk-aware preoperational planning}, while recurrent learning and online
recourse provide adaptive mechanisms whose incremental value is more
instance- and realization-dependent.





% ============================================================
\section{Discussion}
\label{sec:discussion}
% ============================================================

The results support a differentiated view of routing under persistent
disruptions. Temporal persistence matters, but its operational consequences do
not arise simply because average travel times become larger. Explicit
tail-risk optimization provides the clearest and most reproducible improvement
across operational instances, whereas learned search control and online
recourse provide more conditional benefits. This distinction is important
because the three mechanisms address different stages of the decision process:
the disruption model characterizes the environment, CVaR-oriented routing
shapes the preoperational plan, and DDDAS provides state-dependent recourse
during execution.

% ------------------------------------------------------------
\subsection{Why persistence matters even when average impact is controlled}
\label{subsec:discussion_persistence}
% ------------------------------------------------------------

The persistence experiments show that upper-tail service risk depends on the
\emph{temporal organization} of disruption, not only on its average magnitude.
When the Hawkes branching ratio increases while the stationary mean event rate
is held fixed, disruption activity becomes increasingly clustered: quiet
periods become longer, but disruption episodes become more concentrated.

The stricter mean-impact-matched control reinforces this interpretation.
Even after matching the continuous-time mean environmental travel-time
multiplier across persistence regimes, tail tardiness increases substantially
with $\eta$. For the unchanged OR-Tools route, for example,
$\CVaR_{0.95}$ of tardiness rises from approximately $60.7$ to $113.2$
between $\eta=0$ and $\eta=0.70$, even though the route-sampled mean travel-time
multiplier does not increase.

This result illustrates why an average-delay representation can be insufficient
for operational reliability analysis. Two disruption environments with similar
mean impact may generate very different service outcomes if one distributes
delay diffusely over time and the other concentrates it into bursts that arrive
when route slack is already limited.

The mechanism should not be interpreted as feedback from the route into the
Hawkes process. The disruption process remains exogenous. What changes is the
interaction between a temporally clustered environment and the schedule
recursion of the route. Once an early delay changes subsequent departure times,
later arcs may encounter the same external disruption trajectory at less
favorable points in the schedule. Persistence and schedule propagation can
therefore reinforce one another operationally without constituting a
route-to-environment feedback loop.

% ------------------------------------------------------------
\subsection{Tail-risk optimization is more robust than mean optimization}
\label{subsec:discussion_cvar}
% ------------------------------------------------------------

The most consistent algorithmic result is the advantage of explicitly
optimizing the upper tail of tardiness. ALNS-CVaR improves both the aligned
tail-risk objective and $\CVaR_{0.95}$ tardiness relative to ALNS-Mean on all
nine operational days.

The gain is not free. CVaR-oriented solutions travel somewhat farther on
average. This is precisely the trade-off expected from risk-sensitive routing:
additional spatial cost is accepted in exchange for routes whose schedules are
less exposed to severe service failures under unfavorable stochastic
realizations.

Importantly, the evidence does not support the stronger claim that increasing
persistence necessarily makes CVaR optimization progressively more valuable
relative to mean optimization. Persistence strongly increases tail risk, but
the interaction
\[
G(0.70)-G(0)
\]
is not statistically stable across paired repetitions. Thus, two conclusions
should be separated:

\[
\text{persistence increases tail risk},
\]

which is strongly supported, and

\[
\text{persistence necessarily amplifies the relative CVaR-search advantage},
\]

which is not supported by the present data.

This distinction is relevant for practice because it suggests that
tail-sensitive planning can be useful even before the environment reaches an
extreme persistence regime. The value of CVaR does not depend on demonstrating
a monotone interaction with the Hawkes branching ratio; it arises from
protecting service performance against adverse realizations across heterogeneous
operating days.

% ------------------------------------------------------------
\subsection{Learning does not automatically improve adaptive search}
\label{subsec:discussion_rl}
% ------------------------------------------------------------

The learning experiments provide a useful counterpoint to the risk-optimization
results. Neither the feed-forward nor recurrent policy demonstrates a robust
incremental advantage over genuine adaptive ALNS, and neither consistently
outperforms a uniform operator-selection control under the matched online-search
budget.

This result matters because temporal persistence can make recurrent learning
architectures appear conceptually attractive. In the present setting, however,
the observable state already contains stochastic tardiness summaries, capacity
utilization, and route-dependent OMI exposure. A recurrent hidden state may
therefore add limited information beyond these explicit state variables.

The zero-shot transfer experiment further indicates that operator-selection
patterns learned on one operational day are not automatically portable to
other days. This does not imply that reinforcement learning is unsuitable for
vehicle-routing search. Rather, it shows that the added complexity of learning
requires empirical justification under controlled baselines, especially when
offline training cost is substantial.

This qualification is particularly important in applied routing systems.
A learned component should not be regarded as operationally superior merely
because it is more sophisticated than an adaptive heuristic. If a trained policy
does not provide stable gains beyond simpler controls, then the additional
training, maintenance, and validation burden becomes part of the deployment
decision.

% ------------------------------------------------------------
\subsection{Sequential recourse: conditional improvement versus realized outcome}
\label{subsec:discussion_dddas}
% ------------------------------------------------------------

The DDDAS experiment reveals a second distinction that is central to
risk-sensitive online decision making.

Every executed route modification improves the local conditional search
objective on the decision-time scenario bank. Nevertheless, among complete runs
containing at least one route change, realized composite cost improves in only a
minority of cases.

This apparent tension is a natural consequence of stochastic optimization.
At a decision time $t_0$, the controller compares residual routes over a
conditional distribution of possible futures. The selected route is preferable
with respect to this distributional criterion. Execution then reveals only one
future trajectory. A decision that reduces conditional tail risk may therefore
incur additional distance in a realization where the adverse future against
which it provided protection does not occur.

In this sense, online CVaR recourse resembles an insurance decision more than a
deterministic shortest-path correction. The decision may be rational ex ante
even when its insurance premium---typically additional distance---is not repaid
in the subsequently realized trajectory.

The Day~7 results illustrate this mechanism particularly clearly. The mean
realized objective difference is slightly favorable to DDDAS, but the median is
slightly unfavorable and the mean is sensitive to a rare high-consequence
trajectory in which substantial tardiness is avoided. Such events are precisely
the type of outcome that tail-sensitive control is designed to protect against,
but they also imply that average realized gains can be highly variable in
finite samples.

The operational logs reinforce the selective nature of the controller.
Only $16.6\%$ of recourse attempts result in a route modification, and
approximately $37.1\%$ of attempts reach the bounded operator-decision limit.
Moreover, most executed changes do not alter the immediately following
customer. The DDDAS layer should therefore be interpreted as a mechanism for
selective residual-plan reconsideration rather than continual reactive
rerouting.

% ------------------------------------------------------------
\subsection{Managerial implications}
\label{subsec:managerial}
% ------------------------------------------------------------

Several managerial implications follow from these results.

First, route planning under persistent disruption should not be based solely on
average travel-time inflation. The persistence experiments show that two
operating environments with comparable mean travel-time impact can generate
substantially different tail-risk profiles. Managers concerned with service
reliability should therefore monitor not only the magnitude of congestion or
delay, but also its temporal concentration and persistence.

Second, the strongest operational evidence favors incorporating tail risk into
the \emph{preoperational} routing decision. The nine-day comparison shows that
CVaR-oriented ALNS consistently sacrifices a limited amount of distance to
reduce severe tardiness. For delivery systems in which late service carries
disproportionate contractual, reputational, or downstream operational cost,
this trade-off may be preferable to minimizing expected tardiness or route
length alone.

Third, online replanning should not be interpreted as an obligation to change a
route whenever conditions deteriorate. In the confirmatory DDDAS experiment,
most triggers do not lead to a route modification. This is operationally
desirable: unnecessary rerouting can create driver confusion, invalidate
customer communications, increase dispatch workload, and consume computational
resources without improving service.

A useful implementation principle is therefore
\[
\text{detect} \rightarrow
\text{evaluate} \rightarrow
\text{change only if justified},
\]
rather than
\[
\text{detect} \rightarrow
\text{reroute automatically}.
\]

The DDDAS architecture supports precisely this distinction.

% ------------------------------------------------------------
\subsection{Implications for human-in-the-loop planning}
\label{subsec:human_loop}
% ------------------------------------------------------------

The selective behavior of the online controller also suggests a natural role
for human-in-the-loop decision support.

In many delivery operations, dispatchers and drivers retain responsibility for
route changes because optimization systems do not observe every operational
constraint. Temporary road restrictions, parking availability, customer
preferences, driver familiarity, and informal local knowledge may not be
represented explicitly in the routing model.

A memory-aware DDDAS layer can therefore be used as a \emph{decision-support
	filter} rather than as a fully autonomous rerouting system. A high-OMI episode
can flag that the current schedule is operating under elevated persistent
disruption, while conditional suffix optimization can quantify whether an
alternative route appears materially preferable.

This creates at least three possible levels of automation:
\begin{enumerate}
	\item \emph{advisory mode}: the system reports elevated disruption state and
	candidate alternatives to a dispatcher;
	
	\item \emph{confirmation mode}: the system proposes a reroute only when the
	conditional risk improvement exceeds a specified operational threshold, but
	a human approves execution;
	
	\item \emph{autonomous mode}: residual routes are modified automatically
	under preauthorized operating rules.
\end{enumerate}

The results favor a cautious progression across these levels.
Because a conditional improvement does not guarantee a lower realized cost,
decision interfaces should communicate the nature of the gain explicitly.
For example, presenting a recommendation as ``lower predicted tail risk'' is
more informative than presenting it as ``better route.''

This distinction may improve operator trust by aligning the system's
recommendation with the uncertainty actually being optimized.

% ------------------------------------------------------------
\subsection{Strategic implications for service-level agreements}
\label{subsec:sla}
% ------------------------------------------------------------

The results also have implications for service-level agreements (SLAs) and
contract design.

Many logistics contracts are framed around thresholds such as delivery within a
time window, maximum allowed delay, or penalties applied once lateness exceeds a
specified level. Such contracts are intrinsically tail-sensitive even when the
underlying routing system is optimized for average cost.

A mean-oriented routing policy can therefore be misaligned with the economic
structure of the service agreement. Two routing plans with similar expected
tardiness may have very different probabilities of severe delay or very
different upper-tail losses.

CVaR provides a natural operational bridge between these two levels.
Although it is not itself an SLA metric, optimizing
$\CVaR_{0.95}$ explicitly penalizes the subset of scenarios in which service
degradation is most severe. This can be particularly relevant when lateness
cost is nonlinear or when a small number of severe failures dominate customer
dissatisfaction.

The nine-day results suggest that this form of protection can be obtained with a
moderate distance premium. The practical question for an operator is therefore
not simply whether the CVaR route is longer, but whether the additional routing
cost is smaller than the expected economic consequence of severe SLA failure.

This interpretation also suggests that the objective weight $\lambda_T$ should
ultimately be linked to business-specific service economics rather than treated
as a universal constant. The present experiments use a fixed common value to
enable controlled algorithmic comparison; deployment would require calibration
against actual penalty structures and service priorities.

% ------------------------------------------------------------
\subsection{Memory-aware decision support as an anticipatory paradigm}
\label{subsec:anticipatory}
% ------------------------------------------------------------

The concept of memory-aware decision support remains useful, but its meaning
should be stated carefully.

In the present model, ``memory'' does not imply that realized vehicle delays
feed back into the disruption process, nor does it imply an intrinsically
infinite-dimensional non-Markovian system. For the exponential Hawkes kernel,
the relevant disruption history is summarized by the current excitation state
$m(t)$.

The anticipatory value comes instead from using that state to distinguish two
environments that may look similar instantaneously but have different recent
histories and therefore different near-term event intensity.

A conventional reactive system asks:
\[
\text{What is the current travel condition?}
\]

A memory-aware system additionally asks:
\[
\begin{gathered}
\text{What does the recent disruption history imply about}\\
\text{the current persistence state?}
\end{gathered}
\]

This shift is modest mathematically but potentially important operationally.
The OMI is not a forecast that a specific future incident will occur.
Rather, it is a compact state variable that summarizes how strongly recent
events contribute to present disruption intensity.

The resulting decision paradigm is therefore \emph{anticipatory but not
	clairvoyant}. It uses observed history to update the distribution of possible
futures while preserving causal information constraints.

This distinction is especially relevant for DDDAS.
The sequential controller does not inspect the realized future trajectory.
At each departure epoch, it reconstructs a conditional future from the observed
event prefix and completed travel times. The controller can therefore adapt to
the evolving disruption state without relying on retrospective information.

% ------------------------------------------------------------
\subsection{Implications for routing-system architecture}
\label{subsec:architecture}
% ------------------------------------------------------------

Taken together, the experiments suggest a layered architecture for
disruption-prone last-mile routing:
\begin{equation}
\boxed{\begin{gathered}
\text{environmental-state estimation}\\[-1pt]
\downarrow\\[-1pt]
\text{risk-aware preoperational planning}\\[-1pt]
\downarrow\\[-1pt]
\text{selective online recourse}
\end{gathered}}
\end{equation}

The first layer estimates the current persistent disruption state.
The second constructs a route that is robust to adverse service outcomes before
vehicles depart.
The third reevaluates selected residual routes only when newly observed
information makes reconsideration worthwhile.

The empirical evidence in this study is strongest for the second layer.
CVaR-oriented planning provides the most reproducible improvement across
instances. The third layer contributes flexibility, but its incremental realized
value is more variable and must be balanced against computational and
organizational cost.

This architecture also helps clarify the potential role of learning.
RL need not replace the optimization layer. It can instead be introduced only
where repeated operational experience demonstrates that learned search control
adds value beyond adaptive heuristics. The present results indicate that such an
advantage should be established empirically rather than assumed.

% ------------------------------------------------------------
\subsection{Limitations and scope of inference}
\label{subsec:limitations}
% ------------------------------------------------------------

The study has several limitations that delimit the interpretation of the
results.

First, the operational instances are derived from real Athens delivery data,
but the disruption process is experimental. The Hawkes parameters are not
estimated from a historical database of Athens traffic incidents. The
persistence experiments should therefore be interpreted as controlled stress
tests rather than empirical estimates of disruption dynamics in Athens.

Second, the tightened time windows and vehicle capacities are experimental
design choices. They provide a common regime in which service reliability can
be compared, but they are not claimed to reproduce the contractual or fleet
configuration of the original operator.

Third, the Hawkes model is univariate and spatially homogeneous at the
environment level. Real urban disruption is likely to depend on location, road
class, time of day, weather, incident type, and network interactions. A
multivariate or spatially structured point process would provide a richer
representation but would also require correspondingly richer calibration data.

Fourth, the DDDAS optimizer performs local recourse on the residual suffix of
the triggered vehicle while preserving customer assignment. It does not solve
a full fleet-wide dynamic reassignment problem. Consequently, the online results
quantify the value of computationally bounded local recourse rather than the
upper bound obtainable from unrestricted dynamic reoptimization.

Fifth, the deep sequential DDDAS experiment uses three prospectively selected
representative days. The 20 realized paths per day characterize stochastic
within-instance variability, but they do not create additional independent
operational instances. Cross-instance conclusions therefore rely primarily on
the nine-day preoperational analysis.

Sixth, the learning results are conditional on the present state representation,
policy architectures, and training protocol. Failure to observe a robust RL
advantage should not be generalized to all learning-assisted vehicle-routing
methods.

Finally, the mean-impact-matched persistence control equalizes the
continuous-time mean environmental multiplier, not the route-sampled multiplier
for every method. Because route departure times differ, small differences in
route-specific exposure remain possible and are part of the operational
interaction between routing and the external disruption trajectory.

These limitations point to several natural extensions, including spatial or
multivariate disruption models, empirical calibration of persistence from
incident or probe-vehicle data, heterogeneous fleets and service agreements,
fleet-wide dynamic reassignment, and deployment studies in which dispatcher
responses to risk-sensitive recommendations can be observed directly.




























% ============================================================
\section{Conclusions}
\label{sec:conclusions}
% ============================================================

This study examined last-mile routing under temporally persistent disruptions
by separating three distinct decision layers: the external disruption process,
risk-aware preoperational planning, and causal online recourse.

The first main result is that disruption persistence matters beyond average
travel-time inflation. When the Hawkes branching ratio is increased while the
stationary mean event rate is controlled, disruption activity becomes more
temporally clustered. Under the stricter control that also matches the
continuous-time mean environmental travel-time multiplier, upper-tail
tardiness still increases substantially. Thus, routing reliability depends not
only on the average magnitude of disruption but also on how disruption is
organized over time.

The second and strongest algorithmic result concerns preoperational risk-aware
planning. Across all nine Athens-derived operational days,
CVaR-oriented ALNS consistently improves both the aligned tail-risk objective
and $\CVaR_{0.95}$ tardiness relative to mean-oriented ALNS. This improvement
comes with a moderate increase in route distance, illustrating a clear
reliability--efficiency trade-off rather than dominance across every individual
performance measure. The results therefore support explicit tail-risk
optimization when severe service delays carry disproportionate operational or
contractual consequences.

By contrast, additional learning complexity does not yield a robust incremental
advantage in the present experiments. Feed-forward and recurrent
operator-selection policies do not consistently outperform genuine adaptive
ALNS or a uniform operator-selection control under matched online-search
budgets, and zero-shot transfer to other operating days does not reveal a stable
generalization benefit. Recurrence is therefore not automatically valuable
simply because the disruption environment contains temporal persistence.

The sequential DDDAS experiment provides a complementary result. Online recourse
is highly selective: most disruption-triggered reconsiderations retain the
incumbent residual route. Every executed route modification improves the local
conditional risk objective at the decision epoch, but this conditional
improvement does not imply lower realized cost on the single trajectory that
subsequently occurs. The DDDAS layer should therefore be interpreted as bounded
risk-sensitive recourse rather than as a mechanism that guarantees ex-post
cost improvement.

Taken together, the findings support a layered view of decision support under
persistent disruption:
\[
\boxed{\begin{gathered}
\text{persistent-state estimation}\\[-1pt]
\downarrow\\[-1pt]
\text{risk-aware preoperational routing}\\[-1pt]
\downarrow\\[-1pt]
\text{selective causal recourse}
\end{gathered}}
\]
Among these layers, the most robust empirical evidence in the present study
supports explicit CVaR-based planning. Learning-assisted search and dynamic
replanning remain potentially useful extensions, but their incremental value is
more dependent on instance structure, available information, and computational
budget.

Future work should focus on empirically calibrated spatial or multivariate
disruption processes, heterogeneous fleets and service-level costs, richer
fleet-wide dynamic reassignment, and deployment settings in which dispatcher or
driver responses can be incorporated explicitly. Such extensions would allow
the memory-aware framework developed here to move from controlled operational
stress testing toward fully data-calibrated dynamic decision support.














% -------------------------------------------------
% REFERENCES
% -------------------------------------------------

\bibliographystyle{elsarticle-harv} 
\bibliography{cvrptw_references_revised}


\end{document}
